

\section{Middle primes IV: the comparison law $\nu$ and its prices $S,A$}
\label{sec:middle-records}

\paragraph{What this section proves and how to read it.}
This section proves (CP4) for middle atoms: it constructs $\nu$
from $\lambda$ by keeping forcing records, replacing non-forcing
records at fixed actual marginals, and deleting the rest, and
bounds symmetric and skew prices $S,A$. A \emph{record} is one
summand (see \cref{subsec:intro:glossary}); $F_0,F_1,F_2$ are fold
records (two-layer, defined below) and $E_0,E_1$ matching records.
The key idea is mirror pairing: adding a required edge together
with its mirror makes the first-order change purely skew. Inputs:
\cref{sec:middle-frame}. Outputs: middle contributions to $S,A$
for \cref{sec:closing}. A first reading should follow the
keep/replace/delete classification and mirror pairing; per-prime
kernel estimates are reference material.
On a middle block the comparison law
keeps the two-layer root and the aligned fold $F_0$, which accept the six
queries, and changes the other two-layer records: the fold records $F_1,F_2$
are replaced by the signed whole-block fold record of parameter $\tlaw$
(\cref{subsec:oct24:pure-local,subsec:pure-bounds}), a combination of mass one
of three positive accepting product laws, each with the exact actual half
marginals of the record it replaces; the matching records $E_0,E_1$ by the
exact-marginal law $K_1$ (\cref{subsec:record:e_side}); the six merge records
by their actual-marginal mirror repairs (\cref{subsec:merge-repair}); the
two-layer constants by strict Paley-edge laws. The folded merges and
constants and the polarization correction $G$ of the first stage are
deleted. Every replacement keeps the actual half marginals at
$\iota$ of the record it replaces, so every record difference annihilates the
constants on both sides of the reserved cut, and every deleted record is
positive. \Cref{lem:middle-comparison-prices} collects the symmetric, skew and
auxiliary prices of a block; they enter $S$, $A$ and the class bounds as sums
over the pieces of the degree charges times local coefficients evaluated at
the piece start. Later sections use more than the prices: the three-arm
leading stars and remainders of the fold changes and the one-arm stars and
degree-four remainders of the matching changes
(\cref{prop:signed-bounds,lem:oct24:matching-star}) are the inputs of the
positive envelopes (\cref{thm:envelopes}); the coefficient
$d_\tlaw$ of the joint Gram matrix for the reference and signed
records (\cref{prop:signed-bounds}(d)) is the input of the matching sides
(\cref{prop:signed-sides}); the negative cap $k_{\rm D}$
(\cref{lem:old-f-law}(e)) enters the completion
(\cref{lem:front12:reference}); and the centred norm $\kappa_F$ of the
reference fold record (\cref{lem:record:old-f-marginal}) is an input of the
variance credits (\cref{sec:variance}).

The idea of the section is a signed pairing of the fold records. The
reference fold record has the two cross edges $x_0\to y_2$ and $y_0\to x_2$
between its layer-zero labels $x$ and its layer-one labels $y$, and the
queries need one more, $x_1\to y_0$. Adding it alone would change the record
at first order in its symmetric and in its skew part. Adding it together with
the mirror edge $x_0\to y_1$ (the reversal of its image $y_1\to x_0$ under the
reserved reflection) makes the first-order change purely skew, a two-arm star
of norm about $\sqrt2\,r$, and pushes the symmetric change to second order:
this is the mirror pairing (\cref{subsec:mirror}). The signed record adds two
further product laws with coefficients $\pm\tlaw/2$, the signed sum being
taken after the products over the primes of the block. The first-order change
stays purely skew and becomes a three-arm star of norm about
$(1+\tlaw^2+(1-\tlaw)^2)^{1/2}r$ (about $1.2732\,r$ at $\tlaw=127/500$), at the
price of a total variation $1+\tlaw$ per block, whatever the number of primes;
\cref{rem:mirror-why} explains the trade. Two inputs about the reference fold
record are exact rather than perturbative: the pair laws of its actual half
marginal (\cref{lem:actual-features}), and a centred bound for its
coincidences, which also lowers its negative cap
(\cref{lem:old-f-law}(e),(f)).

\subsection{One prime: halves, kernels and channels}
\label{subsec:middle-records:notation}

Throughout the section $p\equiv3\pmod4$ is a prime, $p\ge19$ unless a larger
bound is stated (every middle prime is at least $Q_3=19951586>2^{24}$; the
statements below that need $p\ge2^{20}$ or $p\ge2^{24}$ say so), and
\[
 q=\frac1p,\quad r=p^{-1/2},\quad Z=(1-q)(1-2q),\quad h_0=1-14q,\quad
 e_0=\frac{2q}{(1-3q)(1-4q)},\quad k_p=\frac{p-3}4 .
\]
Let $\chi$ be the quadratic character and write $a\to b$ if $\chi(b-a)=1$; as
$\chi(-1)=-1$ this is a tournament. Finite sets carry their uniform
probability and $\E$ is the corresponding average; $H_0$ is the space of
mean-zero functions on $\Fp$, $P$ the Paley operator
$(Ph)(s)=\E_t\chi(t-s)h(t)$, so $P^*=-P$, $P\one=0$ and $P^*P=qI$ on $H_0$, and
$\mathcal G_p$ the group of the maps $t\mapsto\alpha t+\beta$, $\alpha$ a nonzero
square, acting diagonally on tuples of labels (\cref{lem:paley-triple}).

\paragraph{Halves and kernels.}
A \emph{half} is a triple $x=(x_0,x_1,x_2)$ of distinct labels. Let $H_F$ be
the set of halves with $x_0\to x_2$, so $|H_F|=p(p-1)(p-2)/2$, and $\nu_F$ its
uniform law; its coordinates are uniform, $V_i\colon L^2(\Fp)\to L^2(\nu_F)$,
$(V_ih)(x)=h(x_i)$, is an isometry and $V_i^*$ is the conditional expectation
onto $x_i$ (\cref{lem:paley-triple}(d)). At the reserved cut a two-layer fold
record of a middle block is, at one prime, a probability measure on
$H_F\times H_F$: the first half carries the layer-zero labels $x$ of the three
objects and the second the layer-one labels $y$
(\cref{lem:record:fold-identification}(b),(d)). We identify such a measure
with its density $K(x,y)$ relative to $\nu_F\otimes\nu_F$ and with the operator
$Kf(x)=\E_yK(x,y)f(y)$ on $L^2(\nu_F)$; every function on $H_F\times H_F$ is
regarded as such a kernel, and $\|\cdot\|$ is the operator norm on
$L^2(\nu_F)$ unless another space is named. The reserved reflection exchanges
$x_i$ and $y_i$; it maps $K$ to $K^\top(x,y)=K(y,x)$, the adjoint operator.
Put $\operatorname{Sym}K=\frac12(K+K^\top)$,
$\operatorname{Skew}K=\frac12(K-K^\top)$ and $\Pi_\perp=I-\Jproj$. If
$m=g\nu_F$ is a probability measure with $g\ge1-\eta$ for some $\eta\in[0,1)$, the \emph{operator of
$K$ on $L^2(m)$} has the same bilinear form,
$\ip{\varphi}{K_m\psi}_m=\int\varphi(x)K(x,y)\psi(y)\,d\nu_F(x)\,d\nu_F(y)$; by
\cref{lem:record:actual-correlation}(b) its norm is at most $\|K\|/(1-\eta)$,
and its symmetric and skew parts are those of $K$, transported in the same
way. If $K,K'$ are densities of probability measures with both half marginals
$m$, the operator of $K'-K$ on $L^2(m)$ is the operator $\Delta$ of the charging
rule \cref{rem:cp:charging}(ii) for the replacement of $K$ by $K'$; the same
holds when $K'$ is a signed density of mass one with both half marginals $m$.

\paragraph{The Paley channels.}
The cross characters of a pair of halves are
\begin{equation}\label{eq:middle:channels}
 \mathsf A=\chi(y_2-x_0),\qquad\mathsf B=\chi(x_2-y_0),\qquad
 \mathsf C=\chi(y_0-x_1),\qquad\mathsf D=\chi(y_1-x_0);
\end{equation}
on six distinct labels $\mathsf A=1$, $\mathsf B=1$, $\mathsf C=1$ and
$\mathsf D=1$ mean $x_0\to y_2$, $y_0\to x_2$, $x_1\to y_0$ and $x_0\to y_1$.
Let $I_\times$ be the indicator that $x$ and $y$ have no common label. Since
$\chi(-1)=-1$, transposition acts by
\begin{equation}\label{eq:middle:transpose}
 (\mathsf A,\mathsf B,\mathsf C,\mathsf D)^\top
 =(\mathsf B,\mathsf A,-\mathsf D,-\mathsf C),\qquad I_\times^\top=I_\times .
\end{equation}
Kernels below are defined at every pair of halves; a channel whose two labels
coincide vanishes there. The edges of $\mathsf A$ and $\mathsf D$ form a star
at $x_0$ with leaves $y_2,y_1$, and those of $\mathsf B$ and $\mathsf C$ a star
at $y_0$ with leaves $x_2,x_1$; the two stars have no common label. A monomial
in the channels is a \emph{star monomial} if all its edges lie in one star
($\mathsf A$, $\mathsf B$, $\mathsf C$, $\mathsf D$, $\mathsf A\mathsf D$,
$\mathsf B\mathsf C$) and a \emph{nonstar} otherwise; the nine nonstars are
listed in \cref{lem:record:projected-nonstar}. This two-star picture organizes
all the estimates of the fold records: star monomials have norms of order $r$,
and nonstars, which contain two disjoint cross edges, of order $q$.

\paragraph{Features.}
For the operators $A_1=V_2P$, $A_2=V_1P$ and
$(A_3h)(x)=-\E_s[\chi(s-x_2)\chi(s-x_1)h(s)]-q\E h$ of
\cref{lem:paley-triple}(e), the kernels of $A_1V_0^*$, $A_2V_0^*$ and $A_3V_0^*$
are $-\mathsf B$, $\mathsf C$ and $\mathsf B\mathsf C-q$, since they depend on
$y$ only through $y_0$, which is uniform; transposing, $\mathsf A=-V_0A_1^*$
and $\mathsf D=-V_0A_2^*$. Put $A_\star=A_2+A_3$ and, for $0\le\tlaw\le1$,
\[
 A_\tlaw=A_2+\tlaw A_1+(1-\tlaw)A_3,\qquad
 \mathbf a_\tlaw=(\tlaw,1,1-\tlaw),\qquad
 |\mathbf a_\tlaw|^2=1+\tlaw^2+(1-\tlaw)^2,
\]
so that $A_\tlaw=A_{\mathbf a_\tlaw}$ in the notation of
\cref{lem:paley-triple}(e), with the \emph{arm vector} $\mathbf a_\tlaw$, and
$A_0=A_\star$. The actual half marginal $m_0=g_0\nu_F$ of the reference fold
record (\cref{lem:old-f-law}(a)) is translation invariant, so its coordinates
are uniform and $V_i$ is also an isometry $H_0\to L^2(m_0)$; write $V_{i,m}^*$
for its adjoint in $L^2(m_0)$, $\mathcal H_i=V_iH_0$, and $\mathcal K_{12}$ for
the orthogonal complement in $L^2(m_0)$ of $\mathcal H_1+\mathcal H_2$ within
the $m_0$-mean-zero functions of $(x_1,x_2)$. The operator $A_3$ maps $H_0$
into the latter space, because $s\mapsto\E_{m_0}[\chi(s-x_2)\chi(s-x_1)]$ is
translation invariant, hence constant. Put $A_3^\circ=\Pi_{\mathcal K_{12}}A_3$,
$A_\star^\circ=A_2+A_3^\circ$ and
$A^\circ_\tlaw=A_2+\tlaw A_1+(1-\tlaw)A_3^\circ$; the three arms $A_2$,
$\tlaw A_1$ and $(1-\tlaw)A_3^\circ$ of $A^\circ_\tlaw$ map into $\mathcal H_1$,
$\mathcal H_2$ and $\mathcal K_{12}$, and $A^\circ_0=A^\circ_\star$.

\begin{lemma}[Two-feature spectral calculation]\label{lem:two-feature-spectrum}
Let $V\colon H_0\to\mathcal H$ be an isometry and $A\colon H_0\to\mathcal H$
satisfy, for real $a,b,c$,
\[
 V^*A=aI+bP,\qquad A^*A=cI,
\]
where $P^*=-P$ and $P^*P=qI$ on $H_0$, and $r=\sqrt q$. Then
$\delta=c-a^2-b^2q\ge0$, and
\begin{align*}
 \|AV^*+VA^*\|&=|a|+\sqrt{c-b^2q},\\
 \|AV^*-VA^*\|&=|b|r+\sqrt{c-a^2}.
\end{align*}
The spectrum of $AV^*+VA^*$ is contained in
$\{0,a\pm\sqrt{c-b^2q}\}$. Thus one pair of Gram identities determines
both the symmetric reference spectrum and the skew spectrum of a mirror or
signed star, in whichever half-marginal Hilbert space the identities hold.
\end{lemma}

\begin{proof}
Put $M=V^*A$ and $A_\perp=A-VM$. Then $V^*A_\perp=0$ and
$A_\perp^*A_\perp=A^*A-M^*M=\delta I$, proving $\delta\ge0$.
If $\delta>0$, write $A_\perp=\sqrt\delta\,U$, with $U$ an isometry
orthogonal to $V$. Both operators vanish off $VH_0\oplus UH_0$, and on
this sum they are
\[
 AV^*+VA^*=\begin{pmatrix}2aI&\sqrt\delta I\\\sqrt\delta I&0\end{pmatrix},
 \qquad
 AV^*-VA^*=\begin{pmatrix}2bP&-\sqrt\delta I\\\sqrt\delta I&0\end{pmatrix}.
\]
The first matrix has eigenvalues $a\pm\sqrt{a^2+\delta}$.
On a $P$-eigenspace of eigenvalue $i\varepsilon r$, $\varepsilon\in\{\pm1\}$,
the second has eigenvalues
$i(b\varepsilon r\pm\sqrt{b^2q+\delta})$; its norm is
$|b|r+\sqrt{b^2q+\delta}$. Substitution gives the formulas. If $\delta=0$,
then $A=VM$, and the operators on $VH_0$ are $2aI$ and $2bP$;
the same conclusions hold directly.
\end{proof}

\paragraph{The reflected transpose.}
The map $\varsigma(x)=(-x_2,-x_1,-x_0)$ is an involution of $H_F$ preserving
$\nu_F$: the relation $s\to t$ becomes $-t\to-s$ under negation, so
$x_0\to x_2$ is preserved, and $\one[x_0\to x_1]$ and $\one[x_1\to x_2]$ are
exchanged. Put $(\mathcal Uf)(x)=f(\varsigma x)$, a self-adjoint unitary on
$L^2(\nu_F)$, and call $K^\varsigma(x,y)=K(\varsigma y,\varsigma x)$ the
\emph{reflected transpose} of a function $K$ on $H_F\times H_F$. As operators
$K^\varsigma=\mathcal UK^\top\mathcal U$. The map
$(x,y)\mapsto(\varsigma y,\varsigma x)$ preserves $\mathsf A$, $\mathsf B$ and
$I_\times$, and $\mathsf C^\varsigma=\chi(y_1-x_2)$. Only this arithmetic
identification of coordinates is used; no negation invariance of an input
function is assumed.

\subsection{The reference fold record at the reserved cut}\label{subsec:old-f}

Before the fold records $F_1,F_2$ can be replaced, the record
they replace must be understood exactly: the reference's own
two-layer fold record at the reserved cut. At one prime this is
the uniform law on pairs of halves $(x,y)$ with six distinct
labels and the two cross edges $x_0\to y_2$ and $y_0\to x_2$ (the
channels $\mathsf A$ and $\mathsf B$ of
\eqref{eq:middle:channels}). The next lemma collects everything
later sections use about it: how many completions $y$ each half
$x$ has and the resulting normalizing constant and actual half
marginal; an exact spectral description of its leading part; its
centred norms and remainder in the actual marginal; and its
negative cap $k_{\rm D}$ at $\iota$, which enters the completion
inequality $v_{\rm D}\ge k_{\rm D}w_{\rm D}$ of
\cref{lem:front12:reference} and the protection inequality
\eqref{eq:protection}. The exact pair laws of its actual half
marginal follow in \cref{lem:actual-features}.

\begin{lemma}[The reference fold record]\label{lem:old-f-law}
Let $p\ge19$, $\Psi_0=I_\times(1+\mathsf A)(1+\mathsf B)$, $z_0=\E\Psi_0$ (the
average over $\nu_F\otimes\nu_F$) and $K_0=\Psi_0/z_0$. On six distinct labels
$\Psi_0$ is $4$ times the indicator of $x_0\to y_2$ and $y_0\to x_2$, and it
vanishes elsewhere; so $K_0$ is uniform on the pairs of halves with six
distinct labels, $x_0\to y_2$ and $y_0\to x_2$. It is the law at $\iota$ of
each two-layer fold record $\Pi_{F_0},\Pi_{F_1},\Pi_{F_2}$ of the middle
reference at one prime (\cref{lem:record:fold-identification}).
\begin{enumerate}[label=(\alph*)]
\item \emph{Fibres and normalizer.} For $x\in H_F$ put
 $e_{01}=\one[x_0\to x_1]$ and $e_{12}=\one[x_1\to x_2]$. The number of $y$ with
 $\Psi_0(x,y)\ne0$ is $N_0(x)=(p-5)(k_p-1)\bigl(\frac{p-3}2-e_{01}-e_{12}\bigr)$,
 and as $\Psi_0^\top=\Psi_0$ the same formula counts the $x$ for a given $y$.
 Both half marginals of $K_0$ equal $m_0=g_0\nu_F$, where $g_0\circ\varsigma=g_0$,
 \[
  g_0=\frac{(p-2)(p-3-2e_{01}-2e_{12})}{(p-3)(p-4)},\qquad
  \|g_0-1\|_\infty=\eta_F:=\frac{2(p-1)}{(p-3)(p-4)}
  =\frac{2q(1-q)}{(1-3q)(1-4q)}\le e_0,
 \]
 the largest value of $g_0$ is $(p-2)/(p-4)=1+2/(p-4)$, and, with
 $h_0\le z_0\le1$,
 \[
  z_0=\frac{(p-3)(p-4)(p-5)(p-7)}{p(p-1)(p-2)^2}
  =1-14q+\frac{q^2(53-273q+364q^2)}{(1-q)(1-2q)^2}.
 \]
 On a block $\Blk$ of at most $n$ primes $p\ge X$ the density of
 $m_{0,\Blk}=\bigotimes_pm_{0,p}$ relative to $\nu_{F,\Blk}=\bigotimes_p\nu_{F,p}$
 is at most $(1+e_0(X))^n\le e^{ne_0(X)}$, and at most $(1-ne_0(X))^{-1}$ if
 $ne_0(X)<1$, where $e_0(X)$ denotes $e_0$ evaluated at $q=1/X$.
\item \emph{Leading part.} As kernels
 $\Psi_0=\Jproj+\mathsf L+\mathsf A\mathsf B-\mathsf C_{\rm rem}$, with
 $\mathsf L=\mathsf A+\mathsf B$ and
 $\mathsf C_{\rm rem}=(1-I_\times)(1+\mathsf A)(1+\mathsf B)$. On $L^2(\nu_F)$:
 \begin{enumerate}[label=(\roman*)]
 \item $\mathsf L$ annihilates the constants and has its spectrum in
  $\{0\}\cup\{q/(1-q)\pm(q-q^3/(1-q)^2)^{1/2}\}$; hence $\mathsf L\succeq-(r-q)I$
  and $\|\mathsf L\|\le r+q/(1-q)\le\Gamma r$, where $\Gamma=1+\varrho+\varrho'$,
  $\varrho=r\sqrt{1+q}/(1-q)$ and $\varrho'=q(1+\varrho)/(1-2q)$
  (\cref{lem:paley-triple}(d));
 \item $\|\mathsf A\mathsf B\|\le2q/(1-q)\le2q/Z$;
 \item $\mathsf C_{\rm rem}$ is nonnegative and symmetric, its row sums are at
  most $q(17-120q+247q^2)/Z\le17q$, and
  $\|\mathsf C_{\rm rem}\|\le17q\le17q/Z\le36q$.
 \end{enumerate}
\item \emph{Centred norms.} $\|\Pi_\perp K_0\Pi_\perp\|\le h_{\rm old}/h_0$ with
 $h_{\rm old}=\Gamma r+(2/Z+36)q$, and on the $m_0$-centred functions of
 $L^2(m_0)$ the operator of $K_0$ has norm at most
 $\kappa_{\rm old}=(h_{\rm old}/h_0+e_0^2)/(1-e_0)$. Moreover
 $K_0=(\Jproj+\mathsf L)/z_0+R_{\rm old}$ with $\|R_{\rm old}\|\le19q/(Zh_0)$.
\item \emph{Negative cap on the uniform law.}
 $K_0\succeq-r\bigl[1+(19/Z-1)r\bigr]h_0^{-1}I$ on $L^2(\nu_F)$.
\item \emph{Blocks.} Let $\Blk$ be a block of primes $p\ge X\ge2^{24}$ with
 least prime $b$ and $\sum_{p\in\Blk}\log p\le L$, let $n_X\ge\lfloor L/\log
 X\rfloor$ be an integer (the certificate takes $n_X=\lfloor L/\log^-X\rfloor$
 with a rational lower bound $\log^-X\le\log X$), $q_X=1/X$,
 $r_X\ge X^{-1/2}$ and $e_X=e_0(X)$, suppose $n_Xe_X<1$, and put
 $\bar H_X=(1-n_Xe_X)^{-1}$ and
 \[
  k_{\rm D}(X,L)=\bar H_X\bigl[1+(e_o(X)-1)r_X\bigr],
 \]
 where $e_o(X)$ is the constant $e_o$ of \eqref{eq:signed-constants} at $(q_X,r_X)$. The notation $k_{\rm D}(X,L)$ suppresses the choices $n_X,r_X$ fixed here. If $[1+(e_o(X)-1)r_X]r_X<1$, then for every real function
 $\varphi$ on $H_{F,\Blk}=\prod_{p\in\Blk}H_{F,p}$
 \[
  \int\varphi(x)\varphi(y)\bigotimes_{p\in\Blk}K_{0,p}\,
  d\nu_{F,\Blk}(x)\,d\nu_{F,\Blk}(y)\ \ge\ -k_{\rm D}(X,L)\,b^{-1/2}
  \int\varphi^2\,d\nu_{F,\Blk} .
 \]
 Hence every two-layer fold component of the reference of $\Blk$ satisfies
 \[
  \Pi_{F_j}\succeq_\iota-k_{\rm D}(X,L)\,b^{-1/2}\operatorname{Diag}_\iota(\Pi^{(1)}_{F_j}),
 \]
 which is the fold case of \cref{lem:front12:reference}(ii), with its reserved
 cap $k_{\rm D}=k_{\rm D}(X,L)$, on every block of a piece starting at $X$.
\item \emph{Actual centred remainder.} Let $p\ge2^{24}$, let $T_{\rm old}$ be the
 operator of $K_0$ on $L^2(m_0)$, the Markov operator of
 \cref{lem:record:old-f-marginal}(c), let $O=T_{\rm old}-\Jproj$ be its
 centred part, and put
 \[
  O_0=V_0PV^*_{2,m}-V_2PV^*_{0,m}=-\bigl(A_1V^*_{0,m}+V_0A_1^{*_m}\bigr),
 \]
 adjoints in $L^2(m_0)$. Then $O=O_0+R_o$ with $\|R_o\|\le e_oq$, where $e_o$
 is the constant of \eqref{eq:signed-constants} at $(q,r)$; $O_0$ is
 self-adjoint, annihilates the constants, has its spectrum in
 $\{0\}\cup\{q/(1-q)\pm(q-q^3/(1-q)^2)^{1/2}\}$, and $\|O_0\|\le\Gamma_Ar$, with
 $\Gamma_A$ of \cref{lem:actual-features}(d). Hence
 $T_{\rm old}\succeq-\bigl[1+(e_o-1)r\bigr]r\,I$ on $L^2(m_0)$, and the centred
 norm of $T_{\rm old}$ is at most $(\Gamma_A+e_or)r$.
\end{enumerate}
\end{lemma}

\begin{proof}
(a) All counts use \cref{lem:paley-triple}(a). Fix $x$. The label $y_0$ is an
in-neighbour of $x_2$ other than $x_0$, and other than $x_1$ if $e_{12}=1$:
$\frac{p-1}2-1-e_{12}$ choices. Then $y_2$ is a common out-neighbour of $x_0$
and $y_0$ other than $x_2$, and other than $x_1$ if $x_0\to x_1$ and
$y_0\to x_1$; if $e_{01}=1$, exactly $k_p-1$ admissible $y_0$ satisfy
$y_0\to x_1$, namely the common in-neighbours of $x_1,x_2$ other than $x_0$.
Finally $y_1$ is any of the $p-5$ remaining labels. This gives $N_0$. The
ordered pairs of sources $x_0\ne y_0$ with two distinct common out-neighbours
$x_2,y_2$ number $p(p-1)k_p(k_p-1)$, and each extends in $(p-4)(p-5)$ ways;
dividing by $|H_F|$ gives the mean $\bar N_0=(p-3)(p-7)(p-4)(p-5)/(8(p-2))$. As
$\Psi_0=4$ on its support, $\E_yK_0(x,y)=4N_0(x)/(|H_F|z_0)$, and averaging
over $x$ gives $z_0=4\bar N_0/|H_F|$; so each half marginal has the density
$N_0/\bar N_0=g_0$. As $e_{01}+e_{12}\in\{0,1,2\}$, $g_0$ takes the values
$1+2/(p-4)$, $1-2/((p-3)(p-4))$ and $1-2(p-1)/((p-3)(p-4))$, the last being the
farthest from one and the first the largest; $\varsigma$ exchanges $e_{01}$
and $e_{12}$. Multiplying out,
$(1-3q)(1-4q)(1-5q)(1-7q)-(1-14q)(1-q)(1-2q)^2=53q^2-273q^3+364q^4$, positive
because $273^2<4\cdot53\cdot364$, which gives the expansion of $z_0$; and
$(1-q)(1-2q)^2-(1-3q)(1-4q)(1-5q)(1-7q)=q(14-123q+385q^2-420q^3)>0$ for
$q\le1/19$ gives $z_0\le1$. On a block, $g_{0,p}\le1+2/(p-4)\le1+e_0(p)\le
1+e_0(X)$ for $p\ge X$; and $1+y\le e^y$ and $e^{ny}\le(1-ny)^{-1}$ for
$0\le ny<1$.

(b) (i) By the kernel identities of \cref{subsec:middle-records:notation},
$\mathsf L=-(A_1V_0^*+V_0A_1^*)$; it annihilates the constants because $A_1$
does and maps into mean-zero functions. On $H_0$, $A_1^*A_1=qI$ and
$M:=V_0^*A_1=-(I+P)/(p-1)$ (\cref{lem:paley-triple}(e)), so \cref{lem:two-feature-spectrum}, with
$a=b=-q/(1-q)$ and $c=q$, gives the nonzero eigenvalues of $\mathsf L$ as
$q/(1-q)\pm\sqrt{q-q^3/(1-q)^2}$. The
bounds follow from $(q-q^3/(1-q)^2)^{1/2}\le r$ and
$\varrho r=q\sqrt{1+q}/(1-q)\ge q/(1-q)$.
(ii) The kernel $\mathsf A\mathsf B$ depends on $(x_0,x_2)$ and $(y_0,y_2)$
and is a product of two single cross edges; for independent uniform labels it
is a tensor product of two operators of norm $r$, and
\cref{lem:record:projection} with the density $2/(1-q)$ of a directed edge on
both sides gives the bound.
(iii) Both factors of $\mathsf C_{\rm rem}$ are nonnegative, and
$\mathsf C_{\rm rem}^\top=\mathsf C_{\rm rem}$. Fix $x$. Then
$\E_y\mathsf A=\E_y\mathsf B=0$. The pair $(y_0,y_2)$ is a uniform directed
edge, with density $(1+\chi(y_2-y_0)-\one_{y_0=y_2})/(1-q)$ relative to two
independent labels; against $\mathsf A\mathsf B=\chi(y_2-x_0)\chi(x_2-y_0)$ the
three terms contribute $0$, $q\chi(x_2-x_0)=q$ and $-q^2$
(\cref{lem:paley-triple}(b) and $x_0\to x_2$). Hence $\E_y\mathsf A\mathsf B=q$
and $\E_y(1+\mathsf A)(1+\mathsf B)=1+q$, while
$\E_y\Psi_0(x,y)=4N_0(x)/|H_F|\ge(p-5)(p-7)^2/(p(p-1)(p-2))$ by (a). The row sum
of $\mathsf C_{\rm rem}$ at $x$ is the difference,
\[
 \frac{p+1}p-\frac{(p-5)(p-7)^2}{p(p-1)(p-2)}
 =\frac{17p^2-120p+247}{p(p-1)(p-2)}=\frac{q(17-120q+247q^2)}Z ,
\]
at most $17q$ because $17-120q+247q^2\le17(1-3q+2q^2)$ for $q\le69/213$; and
$17/Z\le36$. The Schur test for a nonnegative symmetric kernel bounds the norm
by the largest row sum.

(c) $\Pi_\perp\Jproj\Pi_\perp=0$, so by (b)
$\|\Pi_\perp K_0\Pi_\perp\|\le(\Gamma r+2q/Z+36q)/z_0\le h_{\rm old}/h_0$. The bound
in $L^2(m_0)$ is \cref{lem:record:actual-correlation}(a) with $m_L=m_R=m_0$,
$\eta=e_0\ge\eta_F$ and all functions; and
$R_{\rm old}=(\mathsf A\mathsf B-\mathsf C_{\rm rem})/z_0$.

(d) By (b) and $\Jproj\succeq0$,
$\Psi_0\succeq-(r-q+2q/Z+17q/Z)I=-r[1+(19/Z-1)r]I$; divide by $z_0\ge h_0$.

(e), (f): see the end of \cref{subsec:record:f_remainders}. The proof of (f)
uses \cref{lem:actual-features} and \cref{lem:record:conditional-collision}(d),
and that of (e) uses (f).
\end{proof}

The reference uses the cap $k_{\rm D}$ for its fold type
(\cref{lem:front12:reference}(ii)). Its values $k_{\rm D}(X,585)$ at the $33$
piece starts, all below the recorded F21 block caps of \cref{lem:layer12-caps},
are those of \cref{cert:record:old-f-cap}; at $X=Q_3$ the block
count is $n_X=34$, $\bar H_X<1.0000035$, $1+(e_o-1)r_X<1.0014832$ and
$k_{\rm D}<1.0014866$.

\begin{corollary}[Actual marginal and centred norm of the reference fold record]
\label{lem:record:old-f-marginal}
Let $p\ge19$.
\begin{enumerate}[label=(\alph*)]
\item The actual half marginal is $m_0=g_0\nu_F$, with $g_0$ and
 $\|g_0-1\|_\infty=\eta_F$ as in \cref{lem:old-f-law}(a).
\item The joint law of the two $\iota$-halves of the fold record has the density
 $K_0=\Psi_0/z_0$ relative to $\nu_F\otimes\nu_F$, with
 $z_0=(p-3)(p-4)(p-5)(p-7)/(p(p-1)(p-2)^2)\ge1-14q$.
\item The Markov operator $T_{\rm old}$ of this coupling on $L^2(m_0)$ is
 self-adjoint, $T_{\rm old}\one=\one$, and, provided $(2/Z+36)q<1$ and
 $\eta_F<1$,
 \[
  \|T_{\rm old}-\Jproj\|\le\kappa_F(q,r):=\frac{\dfrac{\bar\Gamma r+(2/Z+36)q}
  {1-(2/Z+36)q}+\eta_F^2}{1-\eta_F},\qquad
  \bar\varrho=\frac{r(1+q)}{1-q},\quad\bar\Gamma=1+\bar\varrho
  +\frac{q(1+\bar\varrho)}{1-2q}.
 \]
\item Let $X$ be a nonempty finite set of middle blocks, $j\in\{0,1,2\}$, and
 $\kappa_X$ the largest value of $\kappa_F(1/p,p^{-1/2})$ at the primes $p$ of
 the blocks of $X$. The Markov operator $T_X$ at $\iota$ of
 $\Pi_X=\bigotimes_{\Blk\in X}\Pi_{F_j,\Blk}$ on $L^2(m_X)$,
 $m_X=\bigotimes_\Blk m_{F_j,\Blk}$ the actual half marginal, satisfies
 $\|T_X-\Jproj\|\le\kappa_X$ and $T_X\preceq(1-\kappa_X)\Jproj+\kappa_XI$; hence, if
 $\kappa_X<1$, with $J_X$ the law with independent $\iota$-halves of law $m_X$,
 \[
  Y[J_X]\ge\frac{Y[\Pi_X]-\kappa_X(1+o(1))C_XY}{1-\kappa_X}.
 \]
\end{enumerate}
\end{corollary}

\begin{proof}
(a) and (b) are \cref{lem:old-f-law}(a). (c) The layer swap preserves the
record, so the coupling is symmetric, and $T_{\rm old}$ is self-adjoint and
fixes the constants. As $\varrho\le\bar\varrho$, $\Gamma\le\bar\Gamma$, and
\cref{lem:old-f-law}(b) gives
$\|\Pi_\perp(\Psi_0-\Jproj)\Pi_\perp\|\le\bar\Gamma r+(2/Z+36)q$; with
$z_0\ge1-14q\ge1-(2/Z+36)q$, $\|\Pi_\perp K_0\Pi_\perp\|$ is at most the inner
fraction, and \cref{lem:record:actual-correlation}(a) with $\eta=\eta_F$ gives
(c). (d) Every $\Pi_{F_j,\Blk}$ has the law $\bigotimes_{p\in\Blk}K_{0,p}$ at $\iota$
(\cref{lem:old-f-law}), so $T_X$ is the tensor product of the prime operators
$T_{\rm old}$, each with $\|T_{\rm old}-\Jproj\|\le\kappa_X$;
\cref{lem:sector}(c) gives $\|T_X-\Jproj\|\le\kappa_X$, and as $T_X$ is
self-adjoint and fixes the constants, $T_X-\Jproj\preceq\kappa_X(I-\Jproj)$. The
energy bound is (P3) of \cref{lem:energy-toolkit} at $\iota$ with the outside
law $\mu$, followed by (E1) for the diagonal copy $\operatorname{Diag}_\iota(m_X)$.
\end{proof}

The certificate evaluates $\kappa_F$ at the piece start with
$\bar\varrho=r(1+q)/(1-q)$, which exceeds the exact norm
$\varrho=r\sqrt{1+q}/(1-q)$ of $V_0^*V_2$, and with $1-(2/Z+36)q$, which is
below $z_0$; both substitutions only increase $\kappa_F$, which is admissible
in $(q,r)$. The bound $\kappa_{\rm old}$ of \cref{lem:old-f-law}(c) bounds the
same norm; the variance credits use $\kappa_F$, and the block bounds of the
signed law the actual centred norm $(\Gamma_A+e_or)r$ of
\cref{lem:old-f-law}(f).

The actual half marginal $m_0$ is not uniform, but its pair laws are exact
and simple, and they give the correlations of the feature spaces without any
density-perturbation allowance.

\begin{lemma}[Exact actual features]\label{lem:actual-features}
Let $p\ge19$, $\mu=2/(p-4)$, and let $m_{ij}$ be the law of $(x_i,x_j)$ under
the actual half marginal $m_0$ of \cref{lem:old-f-law}(a).
\begin{enumerate}[label=(\alph*)]
\item $m_{01}(a,b)=m_{12}(a,b)=\dfrac{1-\mu\chi(b-a)}{p(p-1)}$ for $a\ne b$, and
 $m_{02}$ is uniform on the pairs $a\to b$.
\item On $H_0$, $V^*_{0,m}V_2=(P-qI)/(1-q)$ and
 $V^*_{0,m}V_1=V^*_{1,m}V_2=-(I+\mu pP)/(p-1)$. Hence the correlations of the
 singleton feature spaces in $L^2(m_0)$ are
 \[
  \varrho_{02}=\|V^*_{0,m}V_2\|=\frac{r\sqrt{1+q}}{1-q},\qquad
  \varrho_{01}=\|V^*_{0,m}V_1\|=\|V^*_{1,m}V_2\|
  =\frac q{1-q}\Bigl(1+\frac{4q}{(1-4q)^2}\Bigr)^{1/2};
 \]
 $\varrho_{02}$ is the number $\varrho$ of \cref{lem:paley-triple}(d).
\item On $H_0$,
 \begin{gather*}
  A_3^{*_m}A_3=\frac{q(1-2q-q^2)}{1-q}\,I\preceq qI,\qquad
  A_3^{*_m}V_0=\frac{p+4}{p(p-1)(p-4)}\,I,\\
  V^*_{1,m}A_3=-\frac q{p-1}I+\frac\mu{p-1}P,\qquad
  V^*_{2,m}A_3=-\frac q{p-1}I-\frac\mu{p-1}P .
 \end{gather*}
 Hence, as operators from $H_0$ to $L^2(m_0)$, $\|A_1\|=\|A_2\|=r$,
 $\|A_3^\circ\|\le\|A_3\|\le r$ and $\|\Pi_{\mathcal H_1+\mathcal H_2}A_3\|\le J_Fqr$,
 where
 \[
  J_F=\Bigl(\frac2{1-\varrho_{01}}\Bigr)^{1/2}
  \Bigl(\frac r{1-q}+\frac{2q}{(1-q)(1-4q)}\Bigr).
 \]
\item The synthesis map $(h_0,h_1,h_2,h_3)\mapsto V_0h_0+(A_2h_1+A_1h_2+A_3h_3)/r$
 from $H_0^4$ to $L^2(m_0)$ has squared norm at most
 \[
  \Gamma_A=1+\varrho_{02}+2\varrho_{01}+\frac{qr(1+4q)}{(1-q)(1-4q)}
  +\frac{2qr}{1-q}+\frac{4q^2}{(1-q)(1-4q)} .
 \]
\item Let $\Blk$ be a middle block with least prime $b$ and actual marginal
 $m_{0,\Blk}=\bigotimes_pm_{0,p}$, and let $\mathcal H_i(\Blk)$ and
 $\mathcal K_{12}(\Blk)$ be the spaces of \cref{subsec:middle-records:notation}
 for the whole objects. Let $\mathcal V_0$ be the direct sum over $p\in\Blk$
 of the maps $V_0$ at $p$, an isometry from $\bigoplus_pH_0(\mathbb F_p)$
 into $\mathcal H_0(\Blk)$, and let $\mathcal A_1,\mathcal A_2,\mathcal A_3$ be
 the direct sums of the prime arms $A_2$, $A_1$, $A_3^\circ$. Then
 $\mathcal A_1$, $\mathcal A_2$, $\mathcal A_3$ have ranges in $\mathcal H_1(\Blk)$,
 $\mathcal H_2(\Blk)$, $\mathcal K_{12}(\Blk)$ and norms at most $b^{-1/2}$, and
 for real $c_1,c_2,c_3$ and $\mathcal A=c_1\mathcal A_1+c_2\mathcal A_2+c_3\mathcal A_3$
 the star $\mathcal A\mathcal V_0^*-\mathcal V_0\mathcal A^*$ on $L^2(m_{0,\Blk})$
 is the orthogonal sum over $p\in\Blk$ of the prime stars
 $(c_1A_2+c_2A_1+c_3A_3^\circ)V^*_{0,m}-V_0(c_1A_2+c_2A_1+c_3A_3^\circ)^{*_m}$ on the
 sectors of exact support $\{p\}$.
\end{enumerate}
\end{lemma}

\begin{proof}
(a) By \cref{lem:old-f-law}(a) the weight of a half $x$ under $m_0$ is
proportional to $\one[x_0\to x_2](d-1-e_{01}-e_{12})$, with $d=(p-1)/2$; put
$k=k_p=(d-1)/2$. Fix $(x_1,x_2)=(a,b)$ with $a\ne b$ and put $j=\one[a\to b]$.
The label $x_0$ is an in-neighbour of $b$ other than $a$, in $d-j$ ways, and
$k$ of these are also in-neighbours of $a$ (\cref{lem:paley-triple}(a)); so
the total weight is $(d-j)(d-1-j)-k$, which is $(d-1)(d-\frac12)$ for $j=0$ and
$(d-1)(d-\frac52)$ for $j=1$. Their mean over the two values of $j$ is
$(d-1)(d-\frac32)=(d-1)(p-4)/2$, and the two values are this mean times
$1+\mu$ and $1-\mu$; as $\chi(b-a)=2j-1$, this is the formula for $m_{12}$. The
map $\varsigma$ preserves $m_0$ (\cref{lem:old-f-law}(a)) and sends $(x_0,x_1)$
to $(-x_2,-x_1)$, so $m_{01}(a,b)=m_{12}(-b,-a)$, which is the same formula.
For $x_0\to x_2$, the sum of $d-1-e_{01}-e_{12}$ over $x_1\notin\{x_0,x_2\}$ is
$(p-2)(d-1)-2(d-1)$, the same for every directed pair; so $m_{02}$ is uniform
on the directed pairs.

(b) The first identity is that of \cref{lem:paley-triple}(d), as $m_{02}$ is the
uniform law of directed pairs, as under $\nu_F$. For $f\in H_0$, by (a),
$(V^*_{0,m}V_1f)(a)=\sum_{b\ne a}\frac{1-\mu\chi(b-a)}{p-1}f(b)
=-\frac{f(a)+\mu p(Pf)(a)}{p-1}$, and likewise for $V^*_{1,m}V_2$. On $H_0$ the
operator $P$ is normal with eigenvalues $\pm ir$ (\cref{lem:paley-triple}(c)),
so $\|P-qI\|^2=q+q^2$ and $\|I+\mu pP\|^2=1+\mu^2p$, and
$\mu^2p=4q/(1-4q)^2$.

(c) For $h\in H_0$, $A_3h$ is the function
$-\frac1p\sum_s\chi(s-x_1)\chi(s-x_2)h(s)$ of $(x_1,x_2)$, which is symmetric in
$(x_1,x_2)$; and the density of $m_{12}$ relative to the uniform law of
distinct pairs is $1-\mu\chi(x_2-x_1)$, whose nonconstant part is
antisymmetric. So the Gram $A_3^{*_m}A_3$ is that of the uniform law of
distinct pairs. The pair $(x_1,x_2)$ under $\nu_F$ gives the same Gram,
since its density relative to that law is $1-\chi(x_2-x_1)/(p-2)$, and it is the value
$(1-2q-q^2)/(p-1)\cdot I$ of \cref{lem:paley-triple}(e). Next, by (a) and
\eqref{eq:two-root},
\[
 (V^*_{1,m}A_3h)(a)=-\frac1{p(p-1)}\sum_s\chi(s-a)h(s)
 \sum_{b\ne a}\bigl(1-\mu\chi(b-a)\bigr)\chi(s-b)
 =\frac{-h(a)+\mu p(Ph)(a)}{p(p-1)},
\]
because the inner sum is $-\chi(s-a)-\mu$ for $s\ne a$; conditioning on $x_2$
instead, the inner sum is $-\chi(s-b)+\mu$, which gives $V^*_{2,m}A_3$. For
$A_3^{*_m}V_0$, translation invariance shows that
$\E_{m_0}[\chi(s-x_1)\chi(s-x_2)\mid x_0=a]=\kappa(s-a)$ for a function $\kappa$;
the computation below gives $\kappa(0)=\kappa_0$ and $\kappa(t)=\kappa_1$ for every
$t\ne0$, and then $A_3^{*_m}V_0=-q(\kappa_0-\kappa_1)I$ on $H_0$, since
$\ip{A_3h}{V_0h'}_{m_0}=-\E_a\E_sh'(a)h(s)\kappa(s-a)$. Fix $x_0=0$ and write $b=x_1$,
$c=x_2$; the conditional weight is $\one[0\to c](d-1-e_{0b}-e_{bc})$, of total
$d(d-1)(p-4)$. With $e_{0b}=(1+\chi(b))/2$, $e_{bc}=(1+\chi(c-b))/2$ and
\eqref{eq:two-root}, the sum over $b\notin\{0,c\}$ of $\chi(b)$ times the weight
is $-(p-3)$ for every $c$, so $\kappa_0=-d(p-3)/(d(d-1)(p-4))=-2/(p-4)$. For
$s\ne0$ put $\xi=\chi(s)$, $u=\chi(s-c)$ and $v=\chi(c)$; the sums over
$b\notin\{0,c\}$ of $\chi(s-b)$, $\chi(b)\chi(s-b)$ and $\chi(c-b)\chi(s-b)$ are
$-\xi-u$, $1-vu$ and $p\one_{c=s}-1-v\xi$, so the sum over $b$ of the
weight times $\chi(s-b)$ is $\frac12[(v-(p-5))(\xi+u)-p\one_{c=s}]$. Summing
$u$ times this over $c\ne0$ with the factor $\one[0\to c]=(1+v)/2$, and using
$(1+v)(v-(p-5))=-(p-6)(1+v)$, $\sum_{c\ne0}(1+v)u=1-\xi$ and
$\sum_{c\ne0}(1+v)u^2=p-2-\xi$, gives $-(p-6)(p-3)/4$, the term with
$\one_{c=s}$ vanishing because $u=0$ at $c=s$; so
$\kappa_1=-(p-6)/((p-1)(p-4))$ and $-q(\kappa_0-\kappa_1)=(p+4)/(p(p-1)(p-4))$.
Finally $\|A_i\|=\|P|_{H_0}\|=r$ for $i=1,2$, as $V_1,V_2$ are isometries;
$\|A_3\|^2\le q$ by the first identity. The Gram of $(V_1,V_2)$ on $H_0^2$ is at
least $(1-\varrho_{01})I$ by (b), so
$\|\Pi_{\mathcal H_1+\mathcal H_2}A_3\|\le(1-\varrho_{01})^{-1/2}\sqrt2\max_i
\|V^*_{i,m}A_3\|$, and $\|V^*_{i,m}A_3\|\le q/(p-1)+\mu r/(p-1)=
qr\bigl(r/(1-q)+2q/((1-q)(1-4q))\bigr)$.

(d) The Gram of the four maps is a $4\times4$ block operator on $H_0^4$ with
diagonal blocks of norm at most one (by (c), and as $V_0$, $A_2/r$, $A_1/r$ are
isometries). Its off-diagonal blocks have the norms $\varrho_{01}$ (for $V_0$ and
$A_2/r$, since $P/r$ is unitary on $H_0$), $\varrho_{02}$ ($V_0$ and $A_1/r$),
$\varrho_{01}$ ($A_2/r$ and $A_1/r$), $\|A_3^{*_m}V_0\|/r=qr(1+4q)/((1-q)(1-4q))$ and
$\|V^*_{i,m}A_3\|/r\le qr/(1-q)+2q^2/((1-q)(1-4q))$ for $i=1,2$, by (b) and (c).
Every row sum of these norms is at most $\Gamma_A$, and the Schur test bounds
the norm of the Gram by the largest row sum.

(e) The Markov operators of couplings with exact marginals preserve the
sectors of exact prime support (\cref{lem:sector}(a)), and the stars are
defined sectorwise. A function of the prime-$p$ labels lying in
$\mathcal H_i(p)$ lies in $\mathcal H_i(\Blk)$. One lying in $\mathcal K_{12}(p)$
lies in $\mathcal K_{12}(\Blk)$: the conditional expectation under $m_{0,\Blk}$ of a
function of the whole-object label $x_1$ given the prime-$p$ labels is a
function of $x_{1,p}$ alone, hence lies in $\mathcal H_1(p)$ plus the constants,
which is orthogonal to $\mathcal K_{12}(p)$; likewise for $x_2$. The prime pieces
of $\mathcal V_0$ and of each $\mathcal A_i$ have ranges in the mutually
orthogonal sectors of exact support $\{p\}$, so these maps are isometric,
respectively of norm $\max_p\|A\|\le\max_pr_p\le b^{-1/2}$ by (c), and the star
splits into its prime stars.
\end{proof}

\subsection{Adding an edge together with its mirror: first-order change purely skew}\label{subsec:mirror}

This subsection contains the algebraic idea on which the fold
replacement rests. Recall that $x$ and $y$ are the layer-zero and
layer-one triples of a fold record at one prime. For the record
$F_1$ to accept the six queries, one further cross edge is needed
beyond the two reference edges, namely $x_1\to y_0$
(\cref{cert:record:middle-geometry}); its channel is $\mathsf C$.
Imposing only this edge would change the record already at first
order in $r=p^{-1/2}$ in its symmetric part, which is too
expensive. The reserved reflection sends the channel $\mathsf C$
to $-\mathsf D$, where $\mathsf D$ is the channel of the mirror
edge $x_0\to y_1$. Imposing the mirror edge as well cancels the
first-order symmetric change exactly, leaving a purely skew
first-order change and a symmetric change only at second order.
Part~(d) of the next lemma is the one-parameter signed form of
this pairing used in \cref{subsec:oct24:pure-local}; mirror pairing
itself is the case $\tlaw=0$.

\begin{lemma}[Mirror pairing]\label{lem:mirror}
Let $\Phi=(1+\mathsf A)(1+\mathsf B)\bigl[(1+\mathsf C)(1+\mathsf D)-1\bigr]$ and
$\Phi'=(1+\mathsf A)(1+\mathsf B)\bigl[(1+\mathsf C)(1-\mathsf D)-1\bigr]$,
polynomials in the channels \eqref{eq:middle:channels}, read as kernels.
\begin{enumerate}[label=(\alph*)]
\item $\operatorname{Sym}\Phi=\Phi_{\rm s}:=(1+\mathsf A)(1+\mathsf B)\mathsf C\mathsf D$ and
 $\operatorname{Skew}\Phi=\Phi_{\rm a}:=(1+\mathsf A)(1+\mathsf B)(\mathsf C+\mathsf D)$,
 whereas $\operatorname{Skew}\Phi'=0$ and
 $\operatorname{Sym}\Phi'=(1+\mathsf A)(1+\mathsf B)(\mathsf C-\mathsf D-\mathsf C\mathsf D)$.
\item The four monomials $\mathsf C\mathsf D$, $\mathsf A\mathsf C\mathsf D$,
 $\mathsf B\mathsf C\mathsf D$, $\mathsf A\mathsf B\mathsf C\mathsf D$ of $\Phi_{\rm s}$
 are nonstars: each contains the two-edge matching of $\mathsf C$ ($x_1$--$y_0$)
 and $\mathsf D$ ($x_0$--$y_1$), whose endpoints span only two of the four pairs
 of $\{x_0,x_1\}\times\{y_0,y_1\}$. The skew part splits as
 $\Phi_{\rm a}=L_p+\Phi_{\rm a,ns}$ into the star part
 $L_p=\mathsf C+\mathsf D+\mathsf A\mathsf D+\mathsf B\mathsf C$ and the nonstar part
 $\Phi_{\rm a,ns}=\mathsf A\mathsf C+\mathsf B\mathsf D+\mathsf A\mathsf B\mathsf C
 +\mathsf A\mathsf B\mathsf D$.
\item $L_p=A_\star V_0^*-V_0A_\star^*$ on $L^2(\nu_F)$, with $A_\star=A_2+A_3$; it
 annihilates the constants, has mean-zero range, and
 \[
  \|L_p\|=|\beta|r+(\gamma-\lambda^2)^{1/2}\le l:=r\Bigl[\Bigl(\frac{2+2q^2/(1-2q)}
  {1-q}\Bigr)^{1/2}+\frac q{1-q}\Bigr],
 \]
 where $\lambda=\frac{p+2}{p(p-1)(p-2)}$, $\beta=-\frac1{p-1}$ and
 $\gamma=q+\frac{1-2q-q^2}{p-1}+\frac{2q}{(p-1)(p-2)}$. In particular $l/r\to\sqrt2$
 as $q\to0$, and $l/r$ is admissible in $(q,r)$.
\item \emph{The signed pairing.} For $0\le\tlaw\le1$ let
 \[
  \Phi_\tlaw=(1+\mathsf C)\bigl(1+(1-\tlaw)\mathsf B\bigr)
  \bigl[1+\mathsf D+(1+\tlaw)\mathsf A+(1-\tlaw)\mathsf A\mathsf D\bigr]
  -(1+\mathsf A)(1+\mathsf B),
 \]
 and put $\tilde x_\tlaw=\mathsf D+\tlaw\mathsf A+(1-\tlaw)\mathsf A\mathsf D$ and
 $\tilde y_\tlaw=\mathsf C-\tlaw\mathsf B+(1-\tlaw)\mathsf B\mathsf C$. Then
 $\Phi_0=\Phi$,
 \[
  \operatorname{Sym}\Phi_\tlaw=\tilde x_\tlaw\tilde y_\tlaw,\qquad
  \operatorname{Skew}\Phi_\tlaw=L_{\tlaw,p}+K_\tlaw-K_\tlaw^\top,\qquad
  K_\tlaw=\mathsf A\mathsf C\bigl(1+(1-\tlaw)\mathsf B\bigr),
 \]
 where
 $L_{\tlaw,p}=\tilde x_\tlaw+\tilde y_\tlaw=\mathsf C+\mathsf D+\tlaw(\mathsf A-\mathsf B)
 +(1-\tlaw)(\mathsf A\mathsf D+\mathsf B\mathsf C)$ is, on $L^2(\nu_F)$, the operator
 $A_\tlaw V_0^*-V_0A_\tlaw^*$, with $A_\tlaw=A_2+\tlaw A_1+(1-\tlaw)A_3$ of
 \cref{subsec:middle-records:notation}. The factor $\tilde x_\tlaw$ is a rectangular
 star with centre $x_0$, leaves $y_1,y_2$ and coefficients $(1,\tlaw,1-\tlaw)$, and
 $\tilde y_\tlaw$ one with centre $y_0$, leaves $x_1,x_2$ and coefficients
 $-(1,\tlaw,1-\tlaw)$ (\cref{lem:record:projected-nonstar}(b)); $K_\tlaw$ does not
 depend on $y_1$, and its only factor that depends on $x_1$ is $\mathsf C$.
\end{enumerate}
\end{lemma}

\begin{proof}
(a) By \eqref{eq:middle:transpose}, $(1+\mathsf A)(1+\mathsf B)$ and
$\mathsf C\mathsf D\mapsto(-\mathsf D)(-\mathsf C)$ are symmetric,
$\mathsf C+\mathsf D$ is skew and $\mathsf C-\mathsf D$ is symmetric; and
$(1+\mathsf C)(1\pm\mathsf D)-1=\mathsf C\pm\mathsf D\pm\mathsf C\mathsf D$.
(b) is the expansion of the products; the star monomials of $\Phi_{\rm a}$ are
the singletons $\mathsf C,\mathsf D$ and the products $\mathsf A\mathsf D$ (at $x_0$)
and $\mathsf B\mathsf C$ (at $y_0$).
(c) The kernels of $A_2V_0^*$ and $A_3V_0^*$ are $\mathsf C$ and
$\mathsf B\mathsf C-q$ (\cref{subsec:middle-records:notation}); by
\eqref{eq:middle:transpose} the kernel of the adjoint $V_0A_\star^*$ is
$(\mathsf C+\mathsf B\mathsf C-q)^\top=-\mathsf D-\mathsf A\mathsf D-q$, and the
difference is $L_p$. Each $A_j$ annihilates the constants and maps into
mean-zero functions (\cref{lem:paley-triple}(e)), hence so does $L_p$. On $H_0$,
by \cref{lem:paley-triple}(e), $M:=V_0^*A_\star=\lambda I+\beta P$ and
$A_\star^*A_\star=\gamma I$, because the parts of $A_2^*A_3$ and $A_3^*A_2$ that
contain $P$ cancel ($P^*=-P$). Apply \cref{lem:two-feature-spectrum} with
$(a,b,c)=(\lambda,\beta,\gamma)$ to obtain the stated norm. Finally
$|\beta|r=rq/(1-q)$ and
$\gamma/q=\frac{2-3q-q^2}{1-q}+\frac{2q^2}{(1-q)(1-2q)}\le\frac{2+2q^2/(1-2q)}{1-q}$.

(d) The first factor of $\Phi_\tlaw+(1+\mathsf A)(1+\mathsf B)$ is
$(1+\mathsf C)(1+(1-\tlaw)\mathsf B)=1+\mathsf B+\tilde y_\tlaw$ times
$1+\mathsf D+(1+\tlaw)\mathsf A+(1-\tlaw)\mathsf A\mathsf D=1+\mathsf A+\tilde x_\tlaw$.
The two factors are polynomials in disjoint channels, so
\[
 \Phi_\tlaw=(1+\mathsf B)\tilde x_\tlaw+(1+\mathsf A)\tilde y_\tlaw
 +\tilde x_\tlaw\tilde y_\tlaw .
\]
By \eqref{eq:middle:transpose}, $\tilde x_\tlaw^\top=-\mathsf C+\tlaw\mathsf B
-(1-\tlaw)\mathsf B\mathsf C=-\tilde y_\tlaw$ and $\tilde y_\tlaw^\top=-\tilde x_\tlaw$,
so $\Phi_\tlaw^\top=-(1+\mathsf A)\tilde y_\tlaw-(1+\mathsf B)\tilde x_\tlaw
+\tilde x_\tlaw\tilde y_\tlaw$. Hence $\operatorname{Sym}\Phi_\tlaw=\tilde x_\tlaw\tilde y_\tlaw$
and $\operatorname{Skew}\Phi_\tlaw=(\tilde x_\tlaw+\tilde y_\tlaw)+(\mathsf B\tilde x_\tlaw
+\mathsf A\tilde y_\tlaw)$, and
$\mathsf B\tilde x_\tlaw+\mathsf A\tilde y_\tlaw=\mathsf A\mathsf C(1+(1-\tlaw)\mathsf B)
+\mathsf B\mathsf D(1+(1-\tlaw)\mathsf A)=K_\tlaw-K_\tlaw^\top$, since
$K_\tlaw^\top=\mathsf B(-\mathsf D)(1+(1-\tlaw)\mathsf A)$. At $\tlaw=0$ the first
factor is $(1+\mathsf C)(1+\mathsf B)(1+\mathsf A)(1+\mathsf D)$, so $\Phi_0=\Phi$. The
kernels of $A_2V_0^*$, $A_1V_0^*$ and $A_3V_0^*$ are $\mathsf C$, $-\mathsf B$ and
$\mathsf B\mathsf C-q$, so the kernel of $A_\tlaw V_0^*$ is $\tilde y_\tlaw-(1-\tlaw)q$
and that of its adjoint $V_0A_\tlaw^*$ is $-\tilde x_\tlaw-(1-\tlaw)q$; their
difference is $L_{\tlaw,p}$. With $\mathsf D=\chi(y_1-x_0)$, $\mathsf A=\chi(y_2-x_0)$,
$\mathsf C=-\chi(x_1-y_0)$ and $\mathsf B=\chi(x_2-y_0)$, $\tilde x_\tlaw$ and
$\tilde y_\tlaw$ are the stated rectangular stars, and
$K_\tlaw=\chi(y_2-x_0)\cdot\mathsf C(1+(1-\tlaw)\mathsf B)$ involves the labels
$x_0,x_1,x_2,y_0,y_2$ only.
\end{proof}

\begin{remark}[Why the mirror]\label{rem:mirror-why}
For the orientation $(1+\mathsf C)(1-\mathsf D)$ the whole first-order change
$\mathsf C-\mathsf D$ is symmetric, priced in $S$, which enters the closing
through $e^S$. For the pairing of \cref{lem:mirror} the symmetric change is of
order $q$, and the first-order change is the skew two-arm star $L_p$ of norm
about $\sqrt2\,r$: odd skew terms vanish in the expansion of the reserved cost,
and the even products of such stars are priced by positive envelopes
(\cref{thm:envelopes}), after reflection averaging, at a fraction of their
norms. The signed record keeps this structure and lowers the leading norm from
$\sqrt2\,r$ to about $|\mathbf a_\tlaw|r=(1+\tlaw^2+(1-\tlaw)^2)^{1/2}r$ by two further
product laws with coefficients $\pm\tlaw/2$ (\cref{lem:mirror}(d)). The cost is
a total variation $1+\tlaw$ per block, which enters only the vanishing
remainders of the comparison theorem through the variation factor of
\cref{def:cp4}, and a larger second-order symmetric part ($s_F$ of
\cref{prop:signed-bounds}). Replacements that change the cross edge
$y_0\to x_2$ of the reference record, and the limits of positive mixtures, are
discussed in \cref{subsec:obs:pairing}.
\end{remark}

\subsection{The signed combination over a whole block replacing $F_1,F_2$}\label{subsec:oct24:pure-local}

The comparison records of $F_1$ and $F_2$ replace the reference fold record,
at every prime, by three positive accepting kernels repaired to the actual
half marginal $m_0$, and combine their tensor products over the block with the
signed coefficients $1,\tlaw/2,-\tlaw/2$. The repair precedes tensorization and
the signed sum follows it: every product has exactly the product of the actual
prime marginals, and the sector lemma prices the block difference by its prime
differences (\cref{lem:sector}(b) with its closing clause and (h),
\cref{rem:signed-mixture-sectors}). In this subsection the densities of the
components are taken relative to $m_0\otimes m_0$; the operator of such a
density $P$ on $L^2(m_0)$ is its Markov operator
$T[P]\psi(x)=\int P(x,y)\psi(y)\,dm_0(y)$, which is the operator on $L^2(m_0)$,
in the sense of \cref{subsec:middle-records:notation}, of the
$\nu_F\otimes\nu_F$-density $P\,(g_0\otimes g_0)$.

\begin{lemma}[Positive components with the actual half marginal]
\label{lem:signed-components}
Let $p\ge2^{24}$, $0\le\tlaw\le1$ and $(u,v)\in\{(1,1),(-1,1),(-1,-1)\}$, and on
$H_F\times H_F$ put
\[
 U_{uv}=2\,\one[x_1\to y_0]\bigl(1+(1-\tlaw)\mathsf B\bigr)(1+u\mathsf D)(1+v\mathsf A),
 \qquad z_{uv}=\E_{\nu_F\otimes\nu_F}U_{uv},\qquad
 R_{uv}=\frac{U_{uv}}{z_{uv}\,g_0\otimes g_0},
\]
so that $R_{uv}$ is the density of the probability measure
$U_{uv}\,\nu_F\otimes\nu_F/z_{uv}$ relative to $m_0\otimes m_0$.
\begin{enumerate}[label=(\alph*)]
\item $U_{uv}\ge0$; it vanishes unless $x_1\to y_0$, and it is invariant under
 $\mathcal G_p$ acting diagonally on all six labels.
\item For all $x,y\in H_F$,
 $|\E_{y'\sim\nu_F}U_{uv}(x,y')-1|\le qR_{\rm row}$ and
 $|\E_{x'\sim\nu_F}U_{uv}(x',y)-1|\le qR_{\rm col}$, where
 \begin{gather*}
  R_{\rm row}=(5-3\tlaw)+\frac{2[1+(5-3\tlaw)q]}{1-2q}
  +(2-\tlaw)\Bigl(1+\frac{2q}{1-q}\Bigr),\\
  R_{\rm col}=(5-4\tlaw)+\frac{2[1+(5-4\tlaw)q]}{1-2q}+\frac{1+q}{1-2q};
 \end{gather*}
 both are below $10$, and $|z_{uv}-1|\le\eta_c:=10q$.
\item The marginal defects $\alpha=\one-R_{uv}\one$ and $\beta=\one-R_{uv}^*\one$,
 $R_{uv}$ acting on $L^2(m_0)$, are $\mathcal G_p$-invariant and $m_0$-centred, and
 $|\alpha|,|\beta|\le M_c:=\frac{2\eta_c+(1+\eta_c)e_0}{h_c(1-e_0)}$, $h_c=1-\eta_c$.
 The query coupling $J_C=\frac2{1-q}\one[x_1\to y_0]$, a density relative to
 $m_0\otimes m_0$, has both marginals $m_0$, satisfies $J_C\beta=0$ and
 $J_C^*\alpha=0$, and $\|\Pi_\perp J_C\Pi_\perp\|\le\varrho_{02}$ on $L^2(m_0)$, with
 $\Pi_\perp$ the projection off the constants and $\varrho_{02}$ of
 \cref{lem:actual-features}(b).
\item The density
 \[
  P_{uv}=\frac{R_{uv}+J_C\bigl(\alpha(x)+\beta(y)\bigr)+2M_cJ_C}{1+2M_c}
 \]
 relative to $m_0\otimes m_0$ is that of a $\mathcal G_p$-invariant probability
 measure supported on $x_1\to y_0$, with both half marginals $m_0$, and
 $\|\Pi_\perp(T[P_{uv}]-T[R_{uv}])\Pi_\perp\|\le2M_c\bigl(\|\Pi_\perp T[R_{uv}]\Pi_\perp\|
 +2\varrho_{02}\bigr)$.
\end{enumerate}
\end{lemma}

\begin{proof}
(a) For $a\ne b$, $1+\chi(b-a)$ is $2$ if $a\to b$ and $0$ otherwise; the other
three factors lie in $[0,2]$ for all values $0,\pm1$ of the channels; and the
channels are invariant under the maps $t\mapsto\alpha t+\beta$ with $\alpha$ a
nonzero square, acting on all labels.

(b) \emph{Rows.} Fix $x$, and write $2\,\one[x_1\to y_0]=1+\mathsf C-\delta_{10}$,
$\delta_{10}=\one[y_0=x_1]$ (at $y_0=x_1$ the channel $\mathsf C$ vanishes). Then
$U_{uv}=H(1+u\mathsf D)(1+v\mathsf A)-\delta_{10}\bigl(1+(1-\tlaw)\mathsf B\bigr)
(1+u\mathsf D)(1+v\mathsf A)$, where $H=(1+\mathsf C)(1+(1-\tlaw)\mathsf B)\ge0$ is a
function of $y_0$. Under $\nu_F$ the label $y_0$ is uniform, $y_2$ is a uniform
out-neighbour of $y_0$, and $y_1$ is uniform on the other $p-2$ labels. By
\cref{lem:paley-triple}(b), $\E H=1+(1-\tlaw)\E_c\chi(c-x_1)\chi(x_2-c)=
1+(1-\tlaw)q$. Next $\E[\mathsf A\mid y_0=c]=(p\one_{c=x_0}-1-\chi(c-x_0))/(p-1)$,
so $(p-1)\E[H\mathsf A]=H(x_0)-\E H-\E_c[H(c)\chi(c-x_0)]$. Here $H(x_0)$ is $0$
or $2(2-\tlaw)$, as $\chi(x_2-x_0)=1$, and
$\E_c[H(c)\chi(c-x_0)]=-q+(1-\tlaw)q-(1-\tlaw)\E_c\chi((c-x_0)(c-x_1)(c-x_2))$,
where the last average has modulus at most $2r$ by \cref{thm:weil}(a). Hence
$|\E[H\mathsf A]|\le q(3-2\tlaw+q+2r)/(1-q)\le2(2-\tlaw)q$, the last step because
$q+2r+(4-2\tlaw)q\le1$; so $|\E[H(1+v\mathsf A)]-1|\le(5-3\tlaw)q$. Given
$(y_0,y_2)$, $|\E[\mathsf D\mid y_0,y_2]|=|\chi(y_0-x_0)+\chi(y_2-x_0)|/(p-2)\le
2q/(1-2q)$, and $H(1+v\mathsf A)\ge0$, so the term with $u\mathsf D$ is at most
$\frac{2q}{1-2q}[1+(5-3\tlaw)q]$ in modulus. For the excluded
query coincidence $y_0=x_1$ (the removed query of the preceding
remark), at $y_0=x_1$
the factor $1+(1-\tlaw)\mathsf B$ is at most $2-\tlaw$, and with
$\xi=\chi(x_1-x_0)$ the conditional means given $y_0=x_1$ are
$\E\mathsf A=-(1+\xi)/(p-1)$,
$\E\mathsf D=-\xi/(p-2)+(1+\xi)/((p-1)(p-2))$ and
$\E\mathsf A\mathsf D=-(p-3)/((p-1)(p-2))$; the largest value of
$1+u\E\mathsf D+v\E\mathsf A+uv\E\mathsf A\mathsf D$ over $u,v,\xi$ is $1+2/(p-1)$.
These three bounds add up to $qR_{\rm row}$.

\emph{Columns.} Fix $y$, and integrate $x_1$ first: given $(x_0,x_2)$ it is
uniform on the other $p-2$ labels, so
$|\E[\mathsf C\mid x_0,x_2]|=|\chi(y_0-x_0)+\chi(y_0-x_2)|/(p-2)\le2q/(1-2q)$. The
remaining factor $(1+(1-\tlaw)\mathsf B)G$, $G=(1+u\mathsf D)(1+v\mathsf A)$ a function
of $x_0$, is nonnegative; $\E G=1-uvq$ by \cref{lem:paley-triple}(b), and
$\E[\mathsf B\mid x_0]=(p\one_{x_0=y_0}-1-\chi(x_0-y_0))/(p-1)$ is $1$ at
$x_0=y_0$ and $0$ or $-2/(p-1)$ elsewhere, so
$-2q(1+q)/(1-q)\le\E[\mathsf BG]\le4q$. Hence
$|\E[(1+(1-\tlaw)\mathsf B)G]-1|\le(5-4\tlaw)q$, and the $\mathsf C$-term is at most
$\frac{2q}{1-2q}[1+(5-4\tlaw)q]$. The removed query has column mass
$q\,\E[(1+(1-\tlaw)\chi(x_2-y_0))G(x_0)\mid x_1=y_0]\le q(1+q)/(1-2q)$: given
$x_1=c$ the pair $(x_0,x_2)$ is uniform on the $(p-1)(p-2)/2$ directed pairs
avoiding $c$, the sum over $x_2$ of $1+(1-\tlaw)\chi(x_2-c)$ is at most
$(p-1)/2$ for every $x_0$, and $\sum_{x_0}G(x_0)=p(1-uvq)$.

For $q\le1/100$ the two bounds are at most $9.19$ and $8.18$. Averaging the row
estimate over $x$ gives $|z_{uv}-1|\le qR_{\rm row}<10q$.

(c) $R_{uv}\one(x)=\E_{y\sim\nu_F}U_{uv}(x,y)/(z_{uv}g_0(x))$. By (b) the ratio
$\E_yU_{uv}(x,y)/z_{uv}$ is within $2\eta_c/h_c$ of one, and
$|g_0-1|\le e_0$ (\cref{lem:old-f-law}(a)), so
$|\alpha|\le(e_0+2\eta_c/h_c)/(1-e_0)\le M_c$; likewise for $\beta$. Invariance
holds because $U_{uv}$, $\nu_F$ and $g_0$ are invariant, and
$\int\alpha\,dm_0=1-\iint R_{uv}\,dm_0\,dm_0=0$. The label $y_0$ is uniform under
$m_0$ and, for every $x$, $\Pr(x_1\to y_0)=(1-q)/2$; so $J_C\one=\one$, and likewise
$J_C^*\one=\one$. As $\beta$ is translation invariant and $m_0$-centred,
$\E_{m_0}[\beta\mid y_0]=0$ (\cref{lem:spec:affpair}(d)(iii)), and
$(J_C\beta)(x)=\frac2{1-q}\E_{m_0}\bigl[\one[x_1\to y_0]\E[\beta\mid y_0]\bigr]=0$;
likewise $J_C^*\alpha=0$. Finally, as $2\,\one[a\to c]=1+\chi(c-a)-\one_{c=a}$,
the operator of $J_C$ on $L^2(m_0)$ is $V_1SV^*_{0,m}$ with
$S=(\Jproj+P-qI)/(1-q)$ on $L^2(\Fp)$; so
$\Pi_\perp J_C\Pi_\perp=V_1\frac{P-qI}{1-q}V^*_{0,m}\Pi_\perp$, of norm at most
$\|P-qI\|/(1-q)=\varrho_{02}$ on $H_0$.

(d) This is \cref{cor:markov-completion} with $m_X=m_Y=m_0$, $R=R_{uv}$,
$J=J_C$ and $M=M_c$; the density vanishes off $x_1\to y_0$, where both
$R_{uv}$ and $J_C$ vanish, and it is invariant because $R_{uv}$, $J_C$, $\alpha$
and $\beta$ are.
\end{proof}

\begin{remark}[Coincidences of the two halves]
The kernels $U_{uv}$ are defined at every pair of halves; only the query
coincidence $y_0=x_1$ is excluded. At another coincidence $x_i=y_j$ the
channels of that pair vanish, and $U_{uv}\ge0$ there. Acceptance needs only
$x_1\to y_0$ (\cref{cert:record:middle-geometry}); invariance under
$\mathcal G_p$, hence uniform one-slot marginals and uniform queried pairs,
holds on the whole product; and every bound of \cref{prop:signed-bounds} is a
bound for the operator of the kernel on all pairs of halves (stars and
nonstars for independent labels, compressed by the induced densities; the
removed query charged $k_{\rm col}q$). No later statement assumes distinct
labels in the two halves of a fold record: the middle frame (\cref{def:record:frame}) and the dominating matrices of \cref{sec:envelopes}, the
matching Gram and the mean credit use operators on $L^2(m_0)$; the
restorations use only the reference records and the root; admissibility
(O1) allows invariant signed measures and (O2) applies only the records' own
norms at $\iota$ (\cref{def:admissible}); and \ref{it:record:R4} concerns only
tail primes.
\end{remark}

\begin{definition}[The signed whole-block fold record]\label{def:record:mixture}
Let $p\ge2^{24}$ and $0\le\tlaw\le1$; the construction uses $\tlaw=127/500$.
\begin{enumerate}[label=(\alph*)]
\item \emph{Reference law:} $\Psi_0$, $z_0$, $K_0=\Psi_0/z_0$ and $m_0=g_0\nu_F$ of
 \cref{lem:old-f-law}.
\item \emph{Components:} the densities $P_{++,p}$, $P_{-+,p}$, $P_{--,p}$ of
 \cref{lem:signed-components} at every prime $p$, relative to
 $m_{0,p}\otimes m_{0,p}$.
\item \emph{The record of $F_1$} on a middle block $\Blk$ is the signed measure
 with the density
 \begin{equation}\label{eq:signed-law}
  K_{\tlaw,\Blk}=\bigotimes_{p\in\Blk}P_{++,p}
  +\frac\tlaw2\bigotimes_{p\in\Blk}P_{-+,p}-\frac\tlaw2\bigotimes_{p\in\Blk}P_{--,p}
 \end{equation}
 relative to $m_{0,\Blk}\otimes m_{0,\Blk}$, $m_{0,\Blk}=\bigotimes_pm_{0,p}$: the
 signed sum is taken after the three products. It has mass one, both half
 marginals $m_{0,\Blk}$, and total variation at most $1+\tlaw$, whatever the
 number of primes.
\item \emph{The record of $F_2$} is the reflected transpose $K^\varsigma_{\tlaw,\Blk}$,
 taken componentwise at every prime; its reference law is $K_0^\varsigma=K_0$.
\item Both are placed by $F_1$ and $F_2$ (\cref{lem:record:fold-identification}(b)),
 with the weight $w_{\rm D}r$ of the records they replace. At $\tlaw=0$ the record is
 $\bigotimes_pP_{++,p}$, whose raw kernel $U_{++}$ adds to the reference record
 the query edge and its mirror: it is a mirror pairing in the sense of
 \cref{lem:mirror}.
\end{enumerate}
\end{definition}

For one prime put
$K_{\tlaw,p}=P_{++,p}+\frac\tlaw2P_{-+,p}-\frac\tlaw2P_{--,p}$,
the same combination as \eqref{eq:signed-law} at one prime, and
write $D_{\tlaw,p}=T[K_{\tlaw,p}]-T_{\rm old}$ for the record
difference of $F_1$ at $p$, where $T$ is extended linearly to signed densities
and $T_{\rm old}$ is the operator of $K_0$ on $L^2(m_0)$, that is $T$ of the
$m_0\otimes m_0$-density $K_0/(g_0\otimes g_0)$; the difference $D_{\tlaw,p}$ annihilates
the constants on both sides, and $D^{\rm s}_{\tlaw,p}$, $D^{\rm a}_{\tlaw,p}$ are its
symmetric and skew parts. The leading part of $D^{\rm a}_{\tlaw,p}$ is
\begin{equation}\label{eq:signed-leading}
 L^\circ_{\tlaw,p}=A^\circ_\tlaw V_{0,m}^*-V_0(A^\circ_\tlaw)^{*_m},\qquad
 A^\circ_\tlaw=A_2+\tlaw A_1+(1-\tlaw)A_3^\circ,
\end{equation}
on $L^2(m_0)$, adjoints taken in $L^2(m_0)$: the operator $L_{\tlaw,p}$ of
\cref{lem:mirror}(d) with the adjoints of the actual marginal and with $A_3$
replaced by its part orthogonal to the singleton features. It is a three-arm
star with centre $\mathcal H_0$, arms in $\mathcal H_1$, $\mathcal H_2$ and
$\mathcal K_{12}$, and coefficients $(1,\tlaw,1-\tlaw)$; at $\tlaw=0$ it is the
two-arm star $A^\circ_\star V^*_{0,m}-V_0(A^\circ_\star)^{*_m}$ of the mirror pairing.
For $F_2$ the record difference is $\mathcal UD_{\tlaw,p}^*\mathcal U$, since
$K^\varsigma=\mathcal UK^\top\mathcal U$ componentwise and $\mathcal UK_0\mathcal U=K_0$;
$\mathcal U$ is unitary on $L^2(m_0)$ because $g_0\circ\varsigma=g_0$, and the
symmetric and skew parts of $\mathcal UD_{\tlaw,p}^*\mathcal U$ are
$\mathcal UD^{\rm s}_{\tlaw,p}\mathcal U$ and $-\mathcal UD^{\rm a}_{\tlaw,p}\mathcal U$.

\paragraph{Acceptance.}
By \cref{cert:record:middle-geometry}, the six queries need in the two-layer
$F_1$ exactly the cross edge $x_1\to y_0$ and in $F_2$ exactly $x_2\to y_1$;
$F_0$ accepts. Every component $P_{uv,p}$ is supported on $x_1\to y_0$
(\cref{lem:signed-components}(a),(d)), so every product of components accepts
the query of $F_1$ at every prime of the block; as
$\mathsf C^\varsigma=\chi(y_1-x_2)$, the reflected components accept the query of
$F_2$. Every component is $\mathcal G_p$-invariant, so its one-slot marginals are
uniform and, $\mathcal G_p$ being transitive on the ordered pairs $(a,a')$ with
$a\to a'$, its queried pair is uniform on the strict-square pairs.

\begin{certificate}[Query images of the middle records]\label{cert:record:middle-geometry}
The frozen geometry record
\path{research/frontier_2026_10_17/structural/middle6_geometry.json}, with the
labels $x_i=i$ and $y_i=3+i$, gives the following projections of
the six slot queries of $\mathcal E_6$ onto each record (the
queries as they read on the original labels, before any
compression). In the two-layer records $F_0$ accepts, and the only needed
cross edges are $x_1\to y_0$ in $F_1$, $x_2\to y_1$ in $F_2$ and $x_0\to y_1$ in
$E_0$ and in $E_1$. In the one-layer records $E_0,E_1$ the only needed edge is
$0\to1$; at the chamber cut $\vartheta_1$ for $E_0$ and $\vartheta_2$ for $E_1$,
label $2$ lies on both sides of the cut and labels $0$ and $1$ on one side each.
The two checker programs re-derived these facts independently
(\path{research/record_gate_2026_10_18/U01/compute/geometry_components.py},
\path{research/record_gate_2026_10_18/U03/compute/indep_geometry.py}).
\end{certificate}

\subsection{Nonstar and coincidence remainders}\label{subsec:record:f_remainders}

This subsection collects the nonstar, rectangular-star and coincidence bounds
used by the signed fold records (\cref{prop:signed-bounds}) and by the
reference fold record (\cref{lem:old-f-law}(e),(f)). No estimate of this
subsection is applied to the matching records. A polynomial of degree at most
one in each character attains its maximum modulus on $[-1,1]^4$ at a vertex,
so it is bounded by its maximum on $\{\pm1\}^4$ everywhere, including where a
character vanishes.

\begin{lemma}[Projected nonstars and rectangular stars]\label{lem:record:projected-nonstar}
\begin{enumerate}[label=(\alph*)]
\item For each nonstar monomial $m$, $\|m\|\le\pi_mq$, where
 \[
 \begin{gathered}
  \pi_{\mathsf A\mathsf B}=\frac2{1-q},\qquad
  \pi_{\mathsf A\mathsf C}=\pi_{\mathsf B\mathsf D}=\Bigl(\frac2Z\Bigr)^{1/2},\qquad
  \pi_{\mathsf C\mathsf D}=\frac1{1-2q},\\
  \pi_{\mathsf A\mathsf B\mathsf C}=\pi_{\mathsf A\mathsf B\mathsf D}=\frac2{Z^{1/2}},\qquad
  \pi_{\mathsf A\mathsf C\mathsf D}=\pi_{\mathsf B\mathsf C\mathsf D}=\frac{\sqrt2}{1-2q},\qquad
  \pi_{\mathsf A\mathsf B\mathsf C\mathsf D}=\frac2{1-2q}.
 \end{gathered}
 \]
 Every $\pi_m$ is at most $2/Z$. Consequently
 $\|\Phi_{\rm s}\|\le\mathcal N_{\rm s}:=q(3+2\sqrt2)/(1-2q)$ and
 $\|\Phi_{\rm a,ns}\|\le\mathcal N_{\rm a}:=q\bigl(2(2/Z)^{1/2}+4/Z^{1/2}\bigr)$.
\item \emph{Rectangular stars.} For independent uniform labels $z,u,v$ and real
 $a_1,a_2,a_3$, the \emph{rectangular star}
 $a_1\chi(u-z)+a_2\chi(v-z)+a_3\chi(u-z)\chi(v-z)$, as an operator between
 functions of $z$ and functions of $(u,v)$, has norm at most
 $r(a_1^2+a_2^2+a_3^2)^{1/2}$: its Gram on mean-zero functions of $z$ is
 $q(a_1^2+a_2^2+a_3^2-2qa_3^2)$, and on the constants $a_3^2q(1-q)$. A product of
 two rectangular stars on disjoint sets of labels has norm at most the product
 of their norms, and its restriction to $H_F\times H_F$, as a kernel relative to
 $\nu_F\otimes\nu_F$, has norm at most $2/Z$ times that.
\end{enumerate}
\end{lemma}

\begin{proof}
(a) A nonstar monomial is $m=m_1m_2$, with $m_1$ a nonempty product of characters of
the star at $x_0$ and $m_2$ of the star at $y_0$. For independent uniform labels
the operator of $m$ is the tensor product of those of $m_1$ and $m_2$. A single
edge has norm $r$, and a full star has norm $(q-q^2)^{1/2}$, its Gram operator
being $q(1-2q)I+q^2\Jproj$ by \cref{lem:paley-triple}(b); so monomials with two,
three and four edges have norms $q$, $q(1-q)^{1/2}$ and $q(1-q)$. Each $m$
depends on $x$ through the coordinates $I$ and on $y$ through $J$, where
$(I,J)$ is $(\{0,2\},\{0,2\})$, $(\{0,1\},\{0,2\})$, $(\{0,2\},\{0,1\})$,
$(\{0,1\},\{0,1\})$, $(\{0,1,2\},\{0,2\})$, $(\{0,2\},\{0,1,2\})$,
$(\{0,1\},\{0,1,2\})$, $(\{0,1,2\},\{0,1\})$ and $(\{0,1,2\},\{0,1,2\})$ in the order
of the display; \cref{lem:record:projection} multiplies the norms by the square
roots of the two density bounds, which gives the $\pi_m$. The coefficients of
$\Phi_{\rm s}$ are one on $\mathsf C\mathsf D,\mathsf A\mathsf C\mathsf D,
\mathsf B\mathsf C\mathsf D,\mathsf A\mathsf B\mathsf C\mathsf D$, and those of
$\Phi_{\rm a,ns}$ one on $\mathsf A\mathsf C,\mathsf B\mathsf D,\mathsf A\mathsf B
\mathsf C,\mathsf A\mathsf B\mathsf D$; the triangle inequality gives the last
claims.
(b) The Gram kernel of the star, as a function of $(z,z')$, is the average over
$(u,v)$ of the product of the two values; by
$\E_u\chi(u-z)\chi(u-z')=\one_{z=z'}-q$ (\cref{lem:paley-triple}(b)) it is
$(a_1^2+a_2^2)(\one_{z=z'}-q)+a_3^2(\one_{z=z'}-q)^2$, the cross terms containing
an isolated average $\E_u\chi(u-\cdot)=0$; as an operator this is
$q(a_1^2+a_2^2+a_3^2(1-2q))I+q\bigl(-(a_1^2+a_2^2)+a_3^2q\bigr)\Jproj$, with the
stated values on mean-zero functions and on the constants, both at most
$q(a_1^2+a_2^2+a_3^2)$. A product on disjoint labels is a tensor product of
operators. The induced density of $\nu_F$ relative to three independent labels
is $2/Z$ (\cref{lem:record:projection}(b)), and \cref{lem:record:projection}(a)
on both halves gives the factor $2/Z$.
\end{proof}

\begin{lemma}[Conditional collisions]\label{lem:record:conditional-collision}
Let $g$ be a polynomial in $\mathsf A,\mathsf B,\mathsf C,\mathsf D$ of degree at
most one in each, $M_g$ its maximum modulus on $\{\pm1\}^4$, and $\delta_{ij}$ the
indicator of $x_i=y_j$, $i,j\in\{0,1,2\}$.
\begin{enumerate}[label=(\alph*)]
\item $\|(1-I_\times)g\|\le9M_gq$.
\item Fix $(i,j)$. On $\{x_i=y_j\}$ every character becomes $0$, a fixed sign, a
 function of the free labels of one half only (\emph{one-sided}), or the
 character of a pair $(x_{i'},y_{j'})$ of free labels (\emph{cross}). After like
 terms are collected, at most two cross characters remain in each monomial, and
 they form a single edge, a two-edge star or a two-edge matching. Let
 $c^{ij}_0(g)$, $c^{ij}_r(g)$, $c^{ij}_q(g)$ be the sums of the absolute
 coefficients of the monomials with no cross character, with a single cross
 edge or a two-edge star, and with a two-edge matching. Then $\delta_{ij}g$ is
 the direct sum over the common value $v$ of $q$ times conditional operators
 $T^{ij}_v$ between the conditional laws given $x_i=v$ and given $y_j=v$, and
 $\|T^{ij}_v\|\le\|T^{ij,0}_v\|+(2/Z)(c^{ij}_rr+c^{ij}_qq)$, where $T^{ij,0}_v$,
 the conditional operator of the no-cross part, has norm at most $c^{ij}_0$.
\item With $c_\ast(g)=\sum_{i,j}c^{ij}_\ast(g)$,
 \[
  \|(1-I_\times)g\|\le q\Bigl[c_0(g)+\frac2Z\bigl(c_r(g)r+c_q(g)q\bigr)
  +\frac{36M_gq}{1-q}\Bigr],
 \]
 and $c^{ij}_0(g)$ may be replaced in $c_0(g)$ by any bound for
 $\sup_v\|T^{ij,0}_v\|$.
\item \emph{The reference collision in bundled form.} Let $g=(1+\mathsf A)(1+\mathsf B)$,
 so that $(1-I_\times)g=\mathsf C_{\rm rem}$ (\cref{lem:old-f-law}(b)), and on a half
 put $h=\chi(x_1-x_0)$ and $k=\chi(x_2-x_1)$. The no-cross parts of the nine
 substitutions $\delta_{ij}g$ sum to
 \[
 \begin{aligned}
  &4\delta_{00}+4\delta_{22}+\delta_{02}+\delta_{20}+\delta_{11}\\
  &\quad+\delta_{01}\bigl(1+k(y)\bigr)+\delta_{10}\bigl(1+k(x)\bigr)
  +\delta_{12}\bigl(1+h(x)\bigr)+\delta_{21}\bigl(1+h(y)\bigr)=qVCV^*,
 \end{aligned}
 \]
 where $V=(V_0,V_1,V_2,kV_1,hV_1)\colon L^2(\Fp)^5\to L^2(\nu_F)$ and $C$ is the
 symmetric $5\times5$ matrix with diagonal $(4,1,4,0,0)$ and entries $1$ at the
 positions $01$, $02$, $12$, $03$, $24$ and their transposes. Then
 $\|C\|\le56119/10000$; on $H_0^5$ the map $V$ has squared norm at most
 $\Gamma_5=1+2(r+q)/(1-q)+2q(1+r)/((1-q)(1-2q))$; and the constant-input part
 $q\bigl[15\,\one\otimes\one+\one\otimes(h+k)+(h+k)\otimes\one\bigr]$ of $qVCV^*$ is
 annihilated by centring on both sides. Consequently, with $Q_\nu$ the
 projection off the constants in $L^2(\nu_F)$ and $c^\nu_{\rm col}$, $c_{\rm col}$
 of \eqref{eq:signed-constants},
 $\|Q_\nu\mathsf C_{\rm rem}Q_\nu\|\le c^\nu_{\rm col}q$, and on the $m_0$-centred
 functions of $L^2(m_0)$ the operator of $\mathsf C_{\rm rem}$ has norm at most
 $c_{\rm col}q$.
\end{enumerate}
\end{lemma}

\begin{proof}
(a) Pointwise $|(1-I_\times)g|\le M_g\sum_{i,j}\delta_{ij}$, and for fixed $x$ each
event $y_j=x_i$ has probability $q$, since $y_j$ is uniform; likewise for fixed
$y$. The Schur test gives the bound.
(b) As $x_i$ and $y_j$ are independent and uniform under $\nu_F\otimes\nu_F$,
$\E[\varphi(x)\delta_{ij}g\,\psi(y)]=q\,\E_v\E[\varphi(x)g\,\psi(y)\mid x_i=y_j=v]$
for real $\varphi,\psi$, which is the direct-sum description; Cauchy--Schwarz in
$v$ bounds the norm by $q\sup_v\|T^{ij}_v\|$. Given $x_i=v$ the two free labels of
$x$ are uniform on a set of $(p-1)(p-2)/2$ pairs, for each $i$ and $v$, so their
density relative to two independent uniform labels is $2/Z$, and likewise for
$y$. The coincidence removes every edge of the two stars at $x_i$ or $y_j$, and
inspection of the nine cases shows that the remaining cross edges are none, a
star or a matching. A monomial without a cross character is a product of
one-sided factors of modulus at most one, so its conditional operator has rank
one and norm at most one. For independent labels a single cross edge, times
one-sided factors, has norm at most $r$, a two-edge star at most
$(q-q^2)^{1/2}\le r$, and a two-edge matching at most $q$; conjugation by the two
square-root densities costs the factor $2/Z$, as in \cref{lem:record:projection}.
Values of the kernel off the support of the conditional laws do not matter, so
fixed signs may be substituted first.
(c) Let $\kappa=\sum_{i,j}\delta_{ij}$. As both halves are injective, simultaneous
coincidences form a partial matching between $\{0,1,2\}$ and $\{0,1,2\}$, and
$0\le\kappa-\one_{\kappa\ge1}\le\binom\kappa2$, the right side being the sum of
$\delta_{ij}\delta_{i'j'}$ over the $18$ unordered pairs of coincidences with
$i\ne i'$, $j\ne j'$. Given $x$, the probability that $y$ has two prescribed
distinct labels at two prescribed positions is at most $(p-2)/|H_F|=2q^2/(1-q)$.
So the kernel $(\kappa-\one_{\kappa\ge1})g$ has row and column sums at most
$36M_gq^2/(1-q)$, and the Schur test bounds its norm. Now
$(1-I_\times)g=\sum_{i,j}\delta_{ij}g-(\kappa-\one_{\kappa\ge1})g$, and (b) bounds
each $\delta_{ij}g$.
(d) \emph{Substitutions.} At $x_0=y_0=v$ and at $x_2=y_2=v$ both characters of
$g$ equal $1$ (for instance $\mathsf A=\chi(y_2-v)=1$ as $y_0\to y_2$), so $g=4$.
At $x_0=y_2$ the character $\mathsf A$ vanishes and $\mathsf B$ stays cross; at
$x_2=y_0$, $\mathsf B$ vanishes and $\mathsf A$ stays cross; at $x_1=y_1$ both stay
cross. At $x_0=y_1$, $\mathsf A=\chi(y_2-y_1)=k(y)$ is one-sided and $\mathsf B$
stays cross; at $x_1=y_0$, $\mathsf B=\chi(x_2-x_1)=k(x)$; at $x_1=y_2$,
$\mathsf A=\chi(x_1-x_0)=h(x)$; at $x_2=y_1$, $\mathsf B=\chi(y_1-y_0)=h(y)$. This
gives the displayed no-cross parts, and leaves twelve single cross characters
counted with the norms $1$ or $2$ of their one-sided multipliers (one each at
$02$ and $20$, two at $11$, two each at $01$, $10$, $12$, $21$) and the two-edge
matching $\mathsf A\mathsf B$ at $11$. As operators on $L^2(\nu_F)$,
$\delta_{ij}=qV_iV_j^*$ (the label $y_j$ is uniform), and the one-sided factors
give $\delta_{01}k(y)=qV_0(kV_1)^*$, $\delta_{10}k(x)=q(kV_1)V_0^*$,
$\delta_{12}h(x)=q(hV_1)V_2^*$ and $\delta_{21}h(y)=qV_2(hV_1)^*$; this is $qVCV^*$.

\emph{The matrix.} $C$ is nonnegative and symmetric. For the positive vector
\[
 w=\Bigl(1,\frac{43367}{10^5},1,\frac{17820}{10^5},\frac{17820}{10^5}\Bigr)
\]
the quotients $(Cw)_a/w_a$ are $561187/10^5$, $243367/43367$, $561187/10^5$,
$10^5/17820$ and $10^5/17820$, all below $56119/10000$, so
$\|C\|\le56119/10000$.

\emph{The synthesis maps.} On $H_0$ every Gram entry $V_a^*V_b$,
$a\ne b$, is a
translation- and dilation-invariant operator, $\alpha I+\beta P$; fixing the two
source labels and summing the third, with $D_0=(p-1)(p-2)$, the ten entries are
\[
\begin{array}{c|cccccccccc}
 ab&01&02&03&04&12&13&14&23&24&34\\\hline
 \alpha&-\frac1{p-1}&-\frac1{p-1}&\frac1{D_0}&\frac1{D_0}
 &-\frac1{p-1}&-\frac1{p-2}&-\frac1{p-2}&\frac1{D_0}&\frac1{D_0}&\frac1{p-2}\\[2pt]
 \beta&-\frac p{D_0}&\frac p{p-1}&\frac p{D_0}&\frac p{p-1}
 &-\frac p{D_0}&0&0&-\frac p{p-1}&-\frac p{D_0}&0
\end{array}
\]
(the entries $01$, $02$, $12$ are those of \cref{lem:paley-triple}(d)). As
$\|\alpha I+\beta P\|=(\alpha^2+q\beta^2)^{1/2}$ on $H_0$, the entries $02$, $04$,
$23$ have norm at most $a_5=(r+q)/(1-q)$ and the others at most
$b_5=q(1+r)/((1-q)(1-2q))\le a_5$; every row has at most two entries of the
first kind, so by the Schur test $\|V|_{H_0^5}\|^2\le1+2a_5+2b_5=\Gamma_5$.

\emph{Centring.} Write $V=V^0+V\Jproj$ with $V^0$ the restriction to $H_0^5$.
Then $VCV^*=V^0C(V^0)^*+\sum_{a,b}C_{ab}\,|U_a\rangle\langle U_b|$ with
$U=(V_a\one)_a=(\one,\one,\one,k,h)$, and $\sum_{a,b}C_{ab}U_a\otimes U_b$ is the
displayed constant-input part. Every $V^0_a$ maps into $\nu_F$-mean-zero
functions (for $kV_1$ and $hV_1$ because $\E[k\mid x_1]$ and $\E[h\mid x_1]$ are
constant), and $Q_\nu(\cdot)Q_\nu$ annihilates the constant-input part: each of
its terms has a constant output or a constant input. Hence $\|Q_\nu(qVCV^*)Q_\nu\|\le q\|C\|\Gamma_5$. The single cross
characters and the matching cost at most $(2/Z)(12r+q)q$ by (b), and
$(\kappa-\one_{\kappa\ge1})g$ at most $36M_gq^2/(1-q)=144q^2/(1-q)$ by the proof of
(c); the certified constant $c^\nu_{\rm col}$ uses $288q/(1-q)$, twice this. So
$\|Q_\nu\mathsf C_{\rm rem}Q_\nu\|\le c^\nu_{\rm col}q$.

\emph{Actual centring.} Let $\varphi,\psi$ be $m_0$-centred with
$\|\varphi\|_{m_0}=\|\psi\|_{m_0}=1$. Then $\|\varphi\|_{\nu_F}\le(1-e_0)^{-1/2}$ and
$|\E_{\nu_F}\varphi|=|\E_{m_0}[\varphi(g_0^{-1}-1)]|\le e_0/(1-e_0)$. Split $\varphi$ and
$\psi$ into their $\nu_F$-centred parts and their $\nu_F$-means; the centred-centred
term is at most $c^\nu_{\rm col}q/(1-e_0)$, and the three terms with a mean are at
most $\|\mathsf C_{\rm rem}\|\le17q$ (\cref{lem:old-f-law}(b)(iii)) times
$2e_0/(1-e_0)^{3/2}+e_0^2/(1-e_0)^2$. This is $c_{\rm col}q$.
\end{proof}

\begin{proof}[Proof of \cref{lem:old-f-law}(e),(f)]
\emph{Transport.} For a kernel $K$ relative to $\nu_F\otimes\nu_F$, the operator
of $K$ on $L^2(m_0)$ is $M_{g_0}^{-1}N_KM_{g_0}^{-1}$, where $M_{g_0}$ is
multiplication by $g_0$ and $N_K\psi(x)=\int K(x,y)\psi(y)\,dm_0(y)$ is its
\emph{natural} operator (\cref{lem:record:actual-correlation}(b)). As
$\|M_{g_0}^{-1}-I\|\le e_0/(1-e_0)$, three rules follow. (T1) The operator has norm
at most $\|K\|/(1-e_0)$. (T2) For every operator $N$ on $L^2(m_0)$,
$\|M_{g_0}^{-1}NM_{g_0}^{-1}-N\|\le D_g\|N\|$, with $D_g=(2e_0+e_0^2)/(1-e_0)^2$ of
\eqref{eq:signed-constants}. (T3) For the constant kernel $\one$,
$\Pi_\perp M_{g_0}^{-1}N_{\one}M_{g_0}^{-1}\Pi_\perp$ is the rank-one operator
$\ip{g_0^{-1}-\one}{\cdot}_{m_0}(g_0^{-1}-\one)$, of norm
$\|g_0^{-1}-\one\|^2_{m_0}\le c_g=e_0^2/(1-e_0)^2$. Moreover, the coordinates of
$m_0$ being uniform, the natural operator of a kernel that depends on $y$ only
through $y_j$ is $AV^*_{j,m}$, where $A$ has the same kernel as an operator from
$L^2(\Fp)$; so $N_{\mathsf L}=V_0PV^*_{2,m}-V_2PV^*_{0,m}=O_0$, and
$N_{\mathsf C}=A_2V^*_{0,m}$, $N_{-\mathsf B}=A_1V^*_{0,m}$,
$N_{\mathsf B\mathsf C-q}=A_3V^*_{0,m}$, the natural operator of a transposed kernel
being the adjoint in $L^2(m_0)$.

(f) \emph{The leading operator.} By \cref{lem:actual-features}(a),(b) the pair
$(x_0,x_2)$ is uniform on the directed pairs under $m_0$, so
$V^*_{0,m}V_2=(P-qI)/(1-q)$ on $H_0$ as under $\nu_F$, and $V_0$, $V_2$ are
isometries of $H_0$ into $L^2(m_0)$. \Cref{lem:two-feature-spectrum} therefore
applies with the same Gram data as in (b)(i) to
$O_0=-(A_1V^*_{0,m}+V_0A_1^{*_m})$: it is self-adjoint,
annihilates the constants and has the stated spectrum. In the synthesis of
\cref{lem:actual-features}(d), $O_0/r$ is the image of a $4\times4$ matrix of
norm one, so $\|O_0\|\le\Gamma_Ar$.

\emph{The remainder.} By \cref{lem:old-f-law}(b),
$\Psi_0=\Jproj+\mathsf L+\mathsf A\mathsf B-\mathsf C_{\rm rem}$, and
$O=\Pi_\perp T_{\rm old}\Pi_\perp$ is $1/z_0$ times the centred operator of
$\Psi_0$ on $L^2(m_0)$. In that operator the constant kernel costs $c_g$ by
(T3); $\mathsf L$ gives $O_0$ up to $D_g\|O_0\|\le D_g\Gamma_Ar$ by (T2);
$\mathsf A\mathsf B$ costs $2q/(Z(1-e_0))$ by (b)(ii) and (T1); and
$\mathsf C_{\rm rem}$ costs $c_{\rm col}q$ by
\cref{lem:record:conditional-collision}(d). So the centred operator of
$\Psi_0$ is $O_0+R'$ with $\|R'\|\le r_oq$, $r_o$ of \eqref{eq:signed-constants},
and its norm is at most $(\Gamma_A+r_or)r$. As $h_0\le z_0\le1$ and
$1-z_0\le14q$ (\cref{lem:old-f-law}(a)), $0\le z_0^{-1}-1\le14q/h_0$, and
$O=O_0+R'+(z_0^{-1}-1)(O_0+R')=O_0+R_o$ with
$\|R_o\|\le r_oq+(14q/h_0)(\Gamma_A+r_or)r=(r_o+n_o)q=e_oq$. Finally
$T_{\rm old}=\Jproj+O\succeq O\succeq O_0-e_oqI$, and the smallest eigenvalue
$q/(1-q)-(q-q^3/(1-q)^2)^{1/2}$ of $O_0$ is at least $q-r$; so
$T_{\rm old}\succeq-(r-q+e_oq)I=-[1+(e_o-1)r]rI$, and
$\|O\|\le\Gamma_Ar+e_oq=(\Gamma_A+e_or)r$.

(e) For $p\in\Blk$ put $c_p=[1+(e_o(p)-1)r_p]r_p$, with $e_o(p)$ the value of $e_o$
at $(1/p,p^{-1/2})$. By (f) the operator $T_p$ of $K_{0,p}$ on $L^2(m_{0,p})$, the
Markov operator of a symmetric coupling of $m_{0,p}$ with itself, is
self-adjoint with spectrum in $[-c_p,1]$. The constant $e_o$ is a
nondecreasing function of $(q,r)$ and at least $56119/10000>1$
(\eqref{eq:signed-constants} and the remark after it), so $c/r=1+(e_o-1)r$ is
nondecreasing, and $c_p\le r_pc_X/r_X\le b^{-1/2}[1+(e_o(X)-1)r_X]$, where
$c_X=[1+(e_o(X)-1)r_X]r_X<1$ by hypothesis. By \cref{lem:sector}(c),
$\bigotimes_pT_p\succeq-\max_pc_p\,I$ on $L^2(m_{0,\Blk})$. Its form at $\varphi$ is the
left side of the inequality, and
$\|\varphi\|^2_{m_{0,\Blk}}\le(1+e_X)^{n_X}\int\varphi^2d\nu_{F,\Blk}\le\bar H_X\int\varphi^2d\nu_{F,\Blk}$
by (a), since $\Blk$ has at most $\lfloor L/\log X\rfloor\le n_X$ primes; as
$\bar H_Xc_X/r_X=k_{\rm D}(X,L)$, this is the inequality. The cut form at
$\iota$ of $\Pi_{F_j}$ at a real half-function is the left side, with $\varphi$ the
half-function read through $F_j$, and that of
$\operatorname{Diag}_\iota(\Pi^{(1)}_{F_j})$ is $\int\varphi^2d\nu_{F,\Blk}$, because the
one-layer fold law is $\nu_{F,\Blk}$
(\cref{lem:record:fold-identification}(b),(d)).
\end{proof}

\subsection{Prime and block bounds of the signed fold records}
\label{subsec:pure-bounds}

The leading skew part of a signed fold difference is the three-arm star
$L^\circ_{\tlaw,p}$ of \eqref{eq:signed-leading}; the next proposition bounds it,
the symmetric part and the remainder, on primes and on blocks. Its constants
are the following expressions in $q$, $r$ and $\tlaw$, with $e_0$, $h_0$, $Z$ and
$|\mathbf a_\tlaw|$ of \cref{subsec:middle-records:notation}, $\varrho_{02}$, $J_F$
and $\Gamma_A$ of \cref{lem:actual-features}, and $\eta_c=10q$, $h_c=1-\eta_c$, $M_c$
of \cref{lem:signed-components}:
\begin{equation}\label{eq:signed-constants}
\begin{aligned}
 &D_g=\frac{2e_0+e_0^2}{(1-e_0)^2},\qquad c_g=\frac{e_0^2}{(1-e_0)^2},\qquad
  n_c=\frac{18/Z+3}{1-e_0}+\frac{(1+2q)c_g}q,\\
 &c_c=\sqrt3\,\Gamma_A+\Bigl(n_c+\frac{D_g}q\sqrt3\,\Gamma_Ar\Bigr)r,\qquad
  E_c=\frac{10c_cr}{h_c}+\frac{2M_c}q\Bigl(\frac{c_cr}{h_c}+2\varrho_{02}\Bigr),\\
 &c_+=\sqrt3\,\Gamma_A+\Bigl(n_c+\frac{D_g}q\sqrt3\,\Gamma_Ar+E_c\Bigr)r,\\
 &e_{\rm sg}=\frac{(18-12\tlaw)/Z+k_{\rm col}}{1-e_0}+\frac{(1+2q)c_g}q
  +\frac{D_g}q\sqrt{3+2\tlaw^2}\,\Gamma_Ar+(1+\tlaw)E_c,\\
 &\Gamma_5=1+\frac{2(r+q)}{1-q}+\frac{2q(1+r)}{(1-q)(1-2q)},\qquad
  c^\nu_{\rm col}=\frac{56119}{10000}\Gamma_5+\frac2Z(12r+q)+\frac{288q}{1-q},\\
 &c_{\rm col}=\frac{c^\nu_{\rm col}}{1-e_0}
  +17\Bigl[\frac{2e_0}{(1-e_0)^{3/2}}+\frac{e_0^2}{(1-e_0)^2}\Bigr],\qquad
  r_o=\frac2{Z(1-e_0)}+c_{\rm col}+\frac{c_g}q+\frac{D_g}q\Gamma_Ar,\\
 &n_o=\frac{14}{h_0}(\Gamma_A+r_or)r,\qquad e_o=r_o+n_o,\qquad
  s_p=\frac{2|\mathbf a_\tlaw|^2/Z+k_{\rm col}}{1-e_0}+c_{\rm col}+(1+\tlaw)E_c+n_o,\\
 &a_p=\frac{2(1+(1-\tlaw)^2)^{1/2}/Z+k_{\rm col}}{1-e_0}
  +\frac{D_g}q|\mathbf a_\tlaw|\Gamma_Ar+(1+\tlaw)E_c+2(1-\tlaw)J_Fr,\\
 &c_{\ge2}=(\Gamma_A+e_or)^2+(1+\tlaw)c_+^2,\qquad
  s_F=\max(s_p,c_{\ge2}),\qquad a_F=\max(a_p,c_{\ge2}),\\
 &\ell_\tlaw=(1+e_0)\biggl[\Bigl(\frac{|\mathbf a_\tlaw|^2+2\tlaw q
  +2(1-\tlaw)^2q^2/(1-2q)}{1-q}\Bigr)^{1/2}+\frac{(1+\tlaw)q}{1-q}\biggr]
  +2(1-\tlaw)J_Fq,\\
 &s_\tlaw=\bigl(4-2\tlaw+2\tlaw^2+2(1+\tlaw^2)^{1/2}\bigr)^{1/2},\qquad
  d_\tlaw=s_\tlaw\Gamma_A+(e_o+e_{\rm sg})r,
\end{aligned}
\end{equation}
where $k_{\rm col}=9/4$ if $\tlaw=0$ or $\tlaw\ge2/(p-1)$, and $k_{\rm col}=3$
otherwise. Each expression is built from $q$, $r$ and positive constants by
sums, products, maxima and square roots, and reciprocals of the positive
decreasing denominators $1-q$, $1-2q$, $1-4q$, $Z$, $h_0$, $h_c$, $1-e_0$ and
$1-\varrho_{01}$, after every remainder carrying $q$ has been written as $r$ times
an admissible coefficient ($e_0/q$, $D_g/q$, $c_g/q$ and $M_c/q$ are themselves
such functions); so for a fixed branch of $k_{\rm col}$ it is a nondecreasing
function of $(q,r)$ (\cref{lem:majorant}), and its value at $(q_X,r_X)$, with
$q_X=1/X$ and $r_X\ge X^{-1/2}$, bounds its value at every prime $p\ge X$. Also
$e_o\ge c_{\rm col}\ge56119/10000$. At $\tlaw=127/500$, $k_{\rm col}=9/4$ for every
$p\ge2^{24}$. In words: $c_cr$ bounds the centred norm of an unnormalized
positive component, $qE_c$ its normalization and positive completion, $c_+r$ the
centred norm of a repaired component; $qe_{\rm sg}$ is the distance of the signed
record from its natural leading operator and $qe_o$ that of the reference
record (\cref{lem:old-f-law}(f)); $s_pq$ and $a_pq$ are the prime symmetric and
skew remainders, $c_{\ge2}/b$ the bound on larger supports, $\ell_\tlaw r$ the
leading star, and $d_\tlaw$ the actual coefficient of the direct Gram of the
signed record; $s_0=\sqrt6$.

\begin{proposition}[Prime and block bounds of the signed fold records]
\label{prop:signed-bounds}
Let $0\le\tlaw\le1$, let $\Blk$ be a middle block with least prime $b\ge2^{24}$,
evaluate the constants of \eqref{eq:signed-constants} at $q=1/b$, $r=b^{-1/2}$,
and let $D_\tlaw=T[K_{\tlaw,\Blk}]-\bigotimes_pT_{{\rm old},p}$ be the record
difference of $F_1$ on $\Blk$, an operator on $L^2(m_{0,\Blk})$, with $T[\cdot]$
taken relative to $m_{0,\Blk}\otimes m_{0,\Blk}$.
\begin{enumerate}[label=(\alph*)]
\item $\|\operatorname{Sym}D_\tlaw\|\le s_F/b$.
\item $\operatorname{Skew}D_\tlaw=L_{\tlaw,\Blk}+R_{\tlaw,\Blk}$, where $L_{\tlaw,\Blk}$ is
 $L^\circ_{\tlaw,p}$ on the sector of exact support $\{p\}$, $p\in\Blk$, and zero on
 the other sectors; $\|L_{\tlaw,\Blk}\|\le\ell_\tlaw b^{-1/2}$ and
 $\|R_{\tlaw,\Blk}\|\le a_F/b$. The operator $L_{\tlaw,\Blk}$ is the three-arm star
 $\mathcal A\mathcal V_0^*-\mathcal V_0\mathcal A^*$ of \cref{lem:actual-features}(e)
 with $(c_1,c_2,c_3)=(1,\tlaw,1-\tlaw)$.
\item \emph{One prime.} At a prime $p\ge2^{24}$, with the constants at
 $(1/p,p^{-1/2})$: $\|D^{\rm s}_{\tlaw,p}\|\le s_pq$,
 $\|D^{\rm a}_{\tlaw,p}-L^\circ_{\tlaw,p}\|\le a_pq$ and $\|L^\circ_{\tlaw,p}\|\le\ell_\tlaw r$;
 every centred component $\Pi_\perp T[P_{uv,p}]\Pi_\perp$ has norm at most $c_+r$, and
 the centred reference $O$ at most $(\Gamma_A+e_or)r$; and the centred signed
 record is
 \[
  \Pi_\perp T[K_{\tlaw,p}]\Pi_\perp=O_0+L^m_{\tlaw,p}+R^{\rm new}_{\tlaw,p},\qquad
  L^m_{\tlaw,p}=A_\tlaw V^*_{0,m}-V_0A_\tlaw^{*_m},\qquad
  \|R^{\rm new}_{\tlaw,p}\|\le e_{\rm sg}q,
 \]
 with $O_0$ of \cref{lem:old-f-law}(f) and $A_\tlaw$ unprojected.
\item \emph{The inputs of the direct Gram.} On every exact support $S\subseteq\Blk$
 with $|S|\ge2$, the restrictions of the centred reference
 $\bigotimes_pT_{\rm old,p}$ and of the centred signed record have norms with sum at
 most $c_{\ge2}/b$; and $c_{\ge2}r<2<s_\tlaw\le d_\tlaw$. With the prime
 decompositions of (c) and \cref{lem:old-f-law}(f) and the frame $\Gamma_A$, these
 are the inputs of the direct common-cut Gram of the signed record
 (\cref{lem:record:direct-gram}, \cref{prop:signed-sides}), whose coefficient before the frame factor $\Gamma_A$ and the
 remainder term in $d_\tlaw$ is $s_\tlaw$ and whose actual coefficient is $d_\tlaw$.
\item For $F_2$ the record difference is $\mathcal UD_\tlaw^*\mathcal U$,
 $\mathcal U=\bigotimes_p\mathcal U_p$, with the same bounds and the leading part
 $-\mathcal UL_{\tlaw,\Blk}\mathcal U$.
\end{enumerate}
\end{proposition}

\begin{proof}
Write $\operatorname{op}(K)=M_{g_0}^{-1}N_KM_{g_0}^{-1}$ for the operator on
$L^2(m_0)$ of a kernel $K$ relative to $\nu_F\otimes\nu_F$, with the transport
rules (T1)--(T3) and the natural operators of the star channels from the
proof of \cref{lem:old-f-law}(e),(f). In the frame of
\cref{lem:actual-features}(d), a natural operator $A_xV^*_{0,m}+V_0A_y^{*_m}$, with
$A_x=A_{\mathbf x}$, $A_y=A_{\mathbf y}$ for vectors $\mathbf x,\mathbf y\in\R^3$ in the
notation of \cref{lem:paley-triple}(e), is $r$ times the image of a matrix of norm
$\max(|\mathbf x|,|\mathbf y|)$, so its norm is at most $\max(|\mathbf x|,|\mathbf y|)\Gamma_Ar$.

\emph{Step 1: the components.} Expand
$U_{uv}=(1+\mathsf C)(1+(1-\tlaw)\mathsf B)(1+u\mathsf D)(1+v\mathsf A)
-\delta_{10}(1+(1-\tlaw)\mathsf B)(1+u\mathsf D)(1+v\mathsf A)$. The star monomials
$\mathsf C,\mathsf B,\mathsf B\mathsf C$ and $\mathsf D,\mathsf A,\mathsf A\mathsf D$ have the
natural operators of the proof of \cref{lem:old-f-law}(e),(f), with the constants
$-q$ of $\mathsf B\mathsf C$ and $\mathsf A\mathsf D$ moved into the constant kernel;
the star part is $A_xV^*_{0,m}+V_0A_y^{*_m}$ with
$\mathbf x=(\tlaw-1,1,1-\tlaw)$ and $\mathbf y=-(v,u,uv)$, so
$\max(|\mathbf x|,|\mathbf y|)=\sqrt3$, and the constant kernel has a coefficient of
modulus at most $1+2q$. The nine nonstar monomials have coefficients of modulus
at most one and norms at most $(2/Z)q$ each
(\cref{lem:record:projected-nonstar}(a)); the removed query has norm at most
$3q$ by the Schur test, its row and column masses being at most
$q(2-\tlaw)(1+2q/(1-q))$ and $q(1+q)/(1-2q)$ (proof of
\cref{lem:signed-components}(b)). By (T1)--(T3) the centred operator of
$U_{uv}$ on $L^2(m_0)$ has norm at most
$\sqrt3\,\Gamma_Ar+D_g\sqrt3\,\Gamma_Ar+n_cq=c_cr$. Dividing by $z_{uv}$, with
$|z_{uv}^{-1}-1|\le\eta_c/h_c$, changes it by at most $(10q/h_c)c_cr$, and the
completion of \cref{lem:signed-components}(d) by at most
$2M_c(c_cr/h_c+2\varrho_{02})$; together $E_cq$. Hence
$\|\Pi_\perp T[P_{uv}]\Pi_\perp\|\le c_cr+E_cq=c_+r$.

\emph{Step 2: the signed record.} With the coefficients $c_{++}=1$,
$c_{-+}=\tlaw/2$, $c_{--}=-\tlaw/2$, $\sum c_{uv}=1$ and $\sum|c_{uv}|=1+\tlaw$, and
by \cref{lem:mirror}(d),
\[
\begin{gathered}
 \sum_{uv}c_{uv}U_{uv}=\Psi^{\rm soft}_\tlaw-\delta_{10}G_\tlaw,\qquad
 \Psi^{\rm soft}_\tlaw=\Phi_\tlaw+(1+\mathsf A)(1+\mathsf B),\\
 G_\tlaw=\bigl(1+(1-\tlaw)\mathsf B\bigr)\bigl[1+\mathsf D+(1+\tlaw)\mathsf A
 +(1-\tlaw)\mathsf A\mathsf D\bigr].
\end{gathered}
\]
The star part of $\Psi^{\rm soft}_\tlaw$ is $A_xV^*_{0,m}+V_0A_y^{*_m}$ with
$\mathbf x=\mathbf a_\tlaw-\mathbf e$ and $\mathbf y=-\mathbf a_\tlaw-\mathbf e$,
$\mathbf e=(1,0,0)$: by \cref{lem:mirror}(d) it is the star part $-\mathbf e$ of
$(1+\mathsf A)(1+\mathsf B)$, whose natural operator is $O_0$, plus $L_{\tlaw,p}$,
whose natural operator is $L^m_{\tlaw,p}$; its constants cancel, and
$\max(|\mathbf x|,|\mathbf y|)=(3+2\tlaw^2)^{1/2}$. The nonstars of
$\Psi^{\rm soft}_\tlaw$ are those of the products of $\{\mathsf C,(1-\tlaw)\mathsf B,
(1-\tlaw)\mathsf B\mathsf C\}$ with $\{\mathsf D,(1+\tlaw)\mathsf A,(1-\tlaw)\mathsf A\mathsf D\}$,
of absolute coefficient sum $(3-2\tlaw)\cdot3=9-6\tlaw$, each of norm at most
$(2/Z)q$. For the removed query, on $\{-1,0,1\}^2$
\[
 \bigl|1+\mathsf D+(1+\tlaw)\mathsf A+(1-\tlaw)\mathsf A\mathsf D\bigr|
 \le1+\tlaw+(1-\tlaw)\mathsf D+\mathsf A+\mathsf A\mathsf D ,
\]
and $0\le1+(1-\tlaw)\mathsf B\le2-\tlaw$. With the conditional means of the
proof of \cref{lem:signed-components}(b), the row average of the right side
given $y_0=x_1$ is at most $1+\tlaw$ when $\tlaw\ge2/(p-1)$, so the row mass is
at most $(2-\tlaw)(1+\tlaw)q\le9q/4$; given $x_1=y_0=c$, the column average of
$(1+(1-\tlaw)\chi(x_2-c))$ times the right side is at most
$(1+\tlaw-q)/(1-2q)$, as in the column estimate there, so the column mass is
below $9q/4$. At $\tlaw=0$ the record is $P_{++}$ alone, whose removed query
has row and column masses below $9q/4$ by \cref{lem:signed-components}(b). For
$0<\tlaw<2/(p-1)$ the row average exceeds $1+\tlaw$ by at most
$2/((p-1)(p-2))$, and both masses are below $3q$. By the Schur test the
removed query has norm at most $k_{\rm col}q$. By (T1)--(T3) and Step~1 for the normalizations and
completions of the three components, weighted by $\sum|c_{uv}|=1+\tlaw$, the
centred signed record is $O_0+L^m_{\tlaw,p}+R^{\rm new}_{\tlaw,p}$ with
$\|R^{\rm new}_{\tlaw,p}\|\le e_{\rm sg}q$. This is the last claim of (c).

\emph{Step 3: the prime difference.} Now
$D_{\tlaw,p}=\Pi_\perp(T[K_{\tlaw,p}]-T_{\rm old})\Pi_\perp$, and
\[
 \sum_{uv}c_{uv}U_{uv}-\Psi_0=\Phi_\tlaw-\delta_{10}G_\tlaw+\mathsf C_{\rm rem},
\]
since $\Psi_0=(1+\mathsf A)(1+\mathsf B)-\mathsf C_{\rm rem}$; this kernel has no constant
term. Hence $D_{\tlaw,p}$ is the centred operator of
$\Phi_\tlaw-\delta_{10}G_\tlaw+\mathsf C_{\rm rem}$, plus the normalizations and
completions of the components (at most $(1+\tlaw)E_cq$), minus $(z_0^{-1}-1)$ times
the centred operator of $\Psi_0$, which is self-adjoint and at most $n_oq$ by
the proof of \cref{lem:old-f-law}(f). \emph{Symmetric part.}
$\operatorname{Sym}\Phi_\tlaw=\tilde x_\tlaw\tilde y_\tlaw$ is a product of two rectangular
stars with $|\mathbf a_\tlaw|^2$ as the sum of the squared coefficients of each, so
its norm on $L^2(\nu_F)$ is at most $(2|\mathbf a_\tlaw|^2/Z)q$
(\cref{lem:record:projected-nonstar}(b)); the removed query costs $k_{\rm col}q$;
both are transported by (T1); and $\mathsf C_{\rm rem}$, symmetric, costs
$c_{\rm col}q$ (\cref{lem:record:conditional-collision}(d)). With the
normalizations and the completions this gives $\|D^{\rm s}_{\tlaw,p}\|\le s_pq$.
\emph{Skew part.} $\operatorname{Skew}\Phi_\tlaw=L_{\tlaw,p}+K_\tlaw-K_\tlaw^\top$. The
operator of $L_{\tlaw,p}$ differs from its natural operator $L^m_{\tlaw,p}$ by at most
$D_g|\mathbf a_\tlaw|\Gamma_Ar$ (T2), and $L^m_{\tlaw,p}-L^\circ_{\tlaw,p}
=(1-\tlaw)\bigl(TV^*_{0,m}-V_0T^{*_m}\bigr)$ with
$T=\Pi_{\mathcal H_1+\mathcal H_2}A_3$, of norm at most $2(1-\tlaw)J_Fqr$
(\cref{lem:actual-features}(c)). On the full label spaces, $K_\tlaw$ ignores the
input label $y_1$ and its output has mean zero in $x_1$; so with $\mathcal H$ the
functions independent of the middle label, $K_\tlaw=\Pi_{\mathcal H^\perp}K_\tlaw\Pi_{\mathcal H}$,
$K_\tlaw-K_\tlaw^\top$ is an off-diagonal block, and
$\|K_\tlaw-K_\tlaw^\top\|=\|K_\tlaw\|\le r\cdot(1+(1-\tlaw)^2)^{1/2}r$, the norms of the
one-edge star $\mathsf A$ and of the rectangular star $\mathsf C(1+(1-\tlaw)\mathsf B)$
(\cref{lem:record:projected-nonstar}(b)); its restriction to $H_F\times H_F$ costs
$2/Z$, and (T1) transports it. The removed query costs $k_{\rm col}q/(1-e_0)$
again. The reference collision $\mathsf C_{\rm rem}$, the normalization of the
reference and the constant centring are self-adjoint and are not charged to the
skew part. With the normalizations and completions this gives
$\|D^{\rm a}_{\tlaw,p}-L^\circ_{\tlaw,p}\|\le a_pq$. \emph{Leading star.} On
$L^2(\nu_F)$, $A_\tlaw^*A_\tlaw$ is the scalar
$\gamma_\tlaw=q[\tlaw^2+1+(1-\tlaw)^2(1-2q-q^2)/(1-q)-2\tlaw q/(1-q)
+2(1-\tlaw)^2q^2/((1-q)(1-2q))]$ by \cref{lem:paley-triple}(e) (the parts with $P$ of
the cross Grams cancel), and $\gamma_\tlaw/q$ is at most the quotient under the
square root in $\ell_\tlaw$; and $V_0^*A_\tlaw=\alpha I+\beta P$ with
$|\beta|=(1+\tlaw)q/(1-q)$. \Cref{lem:two-feature-spectrum} gives
$\|A_\tlaw V_0^*-V_0A_\tlaw^*\|\le|\beta|r+\gamma_\tlaw^{1/2}$ on $L^2(\nu_F)$. The
coordinates of $m_0$ being uniform, $V^*_{0,m}=V_0^*M_{g_0}$ and
$A_\tlaw^{*_m}=A_\tlaw^*M_{g_0}$, so $L^m_{\tlaw,p}$ is unitarily equivalent to
$M_{g_0}^{1/2}(A_\tlaw V_0^*-V_0A_\tlaw^*)M_{g_0}^{1/2}$, of norm at most $1+e_0$ times
that; with the projection cost above, $\|L^\circ_{\tlaw,p}\|\le\ell_\tlaw r$. This
proves (c), the claim on the components being Step~1 and that on the
reference \cref{lem:old-f-law}(f).

\emph{Step 4: blocks.} Every component and the reference are couplings with
both marginals $m_{0,p}$, so their Markov operators are sector-preserving, with
centred norms at most $c_+(p)r_p\le c_+b^{-1/2}$ and
$(\Gamma_A+e_o(p)r_p)r_p\le(\Gamma_A+e_or)b^{-1/2}$, both below one, by (c) and the
monotonicity of the constants. By \cref{rem:signed-mixture-sectors}, $D_\tlaw$
vanishes on the empty sector, is $D_{\tlaw,p}$ on the sector $\{p\}$, and on a
sector $S$ with $|S|\ge2$ has norm at most
$(1+\tlaw)\prod_{p\in S}c_+(p)r_p+\prod_{p\in S}(\Gamma_A+e_o(p)r_p)r_p\le c_{\ge2}/b$,
the further centred factors being contractions; the same bound holds for its
symmetric and skew parts, and all these norms are maxima over the sectors.
The operator $L_{\tlaw,\Blk}$ vanishes on the sectors with $|S|\ge2$ and is the
orthogonal sum of the $L^\circ_{\tlaw,p}$ on the singletons, which is the star of
\cref{lem:actual-features}(e); on the singletons (c) applies with $q_p\le1/b$ and
$r_p\le b^{-1/2}$. This proves (a), (b) and the first claim of (d). For the
second, only the factor $1+\tlaw$ of $c_{\ge2}$ depends on $\tlaw$, so $c_{\ge2}$ is at
most its value at $(2^{-24},2^{-12})$ and $\tlaw=1$, which is
$<8$ by outward evaluation of the formula for $c_{\ge2}$ in
\eqref{eq:signed-constants} (the value is $7.043\ldots$); hence
$c_{\ge2}r<8\cdot2^{-12}<2$; and $s_\tlaw^2\ge4-2\tlaw+2\tlaw^2+2\ge4+3/2$ on
$[0,1]$, while $\Gamma_A\ge1$.

(e) $\mathcal U$ is unitary on $L^2(m_{0,\Blk})$ because $g_0\circ\varsigma=g_0$; the
record difference of $F_2$ is $\mathcal UD_\tlaw^*\mathcal U$ componentwise, whose
symmetric and skew parts are $\mathcal U(\operatorname{Sym}D_\tlaw)\mathcal U$ and
$-\mathcal U(\operatorname{Skew}D_\tlaw)\mathcal U$, and conjugation preserves norms and
the sector decomposition.
\end{proof}

\begin{certificate}[The constants of the signed law and of the matching records]
\label{cert:signed-local}
At each of the $33$ piece starts, the constants below have been
evaluated in outward interval arithmetic and satisfy the gates
listed in this certificate ($R_{\rm row},R_{\rm col}<10$,
$c_+r<1$, $c_{\ge2}r<d_\tlaw$, and the others named). In detail,
the functions \texttt{sharp\_bounds}, \texttt{e\_bounds} and
\texttt{old\_negative} of
\path{research/advance_2026_10_24/local/low_domain.py} evaluate the constants of
\eqref{eq:signed-constants}, the coefficients $s_E(X)$, $a_E(X)$, $r_E(X)$ of
\cref{lem:record:exact-marginal-e}(ii) and \cref{lem:oct24:matching-star}, and
$k_{\rm D}(X,585)$ of \cref{lem:old-f-law}(e), in outward rational interval
arithmetic at $q=1/X$, $r=1/\lfloor\sqrt X\rfloor$ and $\tlaw=127/500$, at each of
the $33$ piece starts; the leading cap is computed there by
\texttt{actual\_leading} of
\path{research/continuation_2026_10_24/local/orientation_local.py}, which is the
closed form $\ell_\tlaw$ of \eqref{eq:signed-constants} at the arm vector
$\mathbf a_\tlaw$. The final scalar evaluation calls them through
\path{research/advance_2026_10_24/global/evaluate_signed.py}, selected by the
low-domain flag of the final design file
\path{research/push_2026_10_25/global/refined_L570_source.json}, from
\path{research/improvements_2026_10_25/global/selected.py}, invoked with the selected push design. It asserts
$R_{\rm row},R_{\rm col}<10$, $c_+r<1$, $c_{\ge2}r<d_\tlaw$ and
the centred-norm bound $\|T_{\rm old}P_0\|\le(\Gamma_A+e_or)r$
of \cref{lem:old-f-law}(f) with a constant exceeding $e_o$, so
that
$(\Gamma_A+e_or)r<1$;
$e_o>1$, $n_Xe_X<1$ and $[1+(e_o-1)r]r<1$; the gates of
\cref{lem:record:exact-marginal-e}(ii); and in each piece the inequalities
$A^{\rm lead}_{F,i}\ge A^{\rm arm}_{F,i}$ and $A_{E,i}\ge A^{\rm arm}_{E,i}$ between
the prices of \cref{subsec:middle-comparison}. The totals of those prices at the
exponent $a$ are stored in \path{research/push_2026_10_25/artifacts/author.json}
(key \texttt{raw\_totals}); no stored record holds the per-piece constants. A
read-only evaluation of the three functions at $X=Q_3=19951586$
($r_X=1/4466$) gives the following upper bounds: $\Gamma_A<1.0002241$,
$R_{\rm row}<7.9840009$, $R_{\rm col}<6.9840008$, $c_+<1.7371505$,
$E_c<0.04070887$, $e_{\rm sg}<17.254638$, $c^\nu_{\rm col}<5.6198028$,
$c_{\rm col}<5.6198068$, $e_o<7.6238442$, $s_F=s_p<11.166062$,
$a_F=a_p<4.7973987$, $c_{\ge2}<4.7880514$, $\ell_\tlaw<1.2731978$,
$s_\tlaw<2.3842273$ and $d_\tlaw<2.3903321$; for the matching records
$s_E<5.4071859\cdot10^{-4}$, $a_E<1.0000003$, $1/(1-q_X)<1.0000001$,
$\kappa^E/r<1.0012128$ and $r_E<10.829703$; and for the cap $n_X=34$,
$\bar H_X<1.0000035$, $1+(e_o-1)r_X<1.0014832$ and $k_{\rm D}<1.0014866$. As
$q\to0$, $\ell_\tlaw\to|\mathbf a_\tlaw|=1.27319\ldots$ and $d_\tlaw\to s_\tlaw$.
Independently, \path{research/advance_2026_10_24/local/check_signed_law.py} checks the
polynomial and leading skew identities of \cref{lem:mirror}(d), the exact
two-centre Gram and $18$ collision sign patterns, with exact finite checks at
small primes; \path{research/advance_2026_10_24/local/check_signed_sharp.py} the
factored symmetric and skew identities, the matrix $C$ of
\cref{lem:record:conditional-collision}(d) with its exact Schur rows, and exact
conditional feature Grams; and
\path{research/advance_2026_10_24/local/check_low_domain.py} the uniform analytic
majorants and gates from $2^{20}$, for $21$ values of the parameter.
\end{certificate}

\subsection{The matching records $E_0,E_1$ and their replacement $K_1$}\label{subsec:record:e_side}

The two two-layer matching records $E_0,E_1$ of a middle block, of weight
$w_{\rm S}r$ each, do not accept the six queries: each needs exactly the cross
edge $x_0\to y_1$ (\cref{cert:record:middle-geometry}). They are replaced by an
exact-marginal law that imposes this edge and its mirror image $x_1\to y_0$, and
nothing else. Unlike the fold records it needs no repair: its fibre count is
constant (\cref{lem:record:e-law}(a)).

\paragraph{The empty half law.}
Let $H_E$ be the set of all halves, $|H_E|=p(p-1)(p-2)$, and $\nu_E$ its uniform
law. Kernels are densities relative to $\nu_E\otimes\nu_E$ and operators act on
$L^2(\nu_E)$; $V_i$, $\Pi_\perp$, $\top$, $\operatorname{Sym}$ and
$\operatorname{Skew}$ are as in \cref{subsec:middle-records:notation}, with $\nu_E$
in place of $\nu_F$. For $i\ne j$ and mean-zero $g$, $\E[g(x_j)\mid x_i]=-g(x_i)/(p-1)$,
so the Gram operator of $V=(V_0,V_1,V_2)$ on $H_0^3$ is $(1+u)I-uJ_3$, with
$u=1/(p-1)$ and $J_3$ the all-ones $3\times3$ matrix; its eigenvalues are
$1+u=1/(1-q)$ and $1-2u=(1-3q)/(1-q)$. Every pair of coordinates of $\nu_E$ is
uniform on the distinct pairs, of density $1/(1-q)$, and $\nu_E$ has density
$1/Z$ (\cref{lem:record:projection}(b)); \cref{lem:record:projection}(a)
applies with these bounds. Write
$\delta_{\mathsf C}=\one[x_1=y_0]$, $\delta_{\mathsf D}=\one[x_0=y_1]$ and
$\delta_{ij}=\one[x_i=y_j]$. At $\iota$ the reference's two-layer matching record
is the uniform law of pairs of disjoint halves, $K^E_0=I_\times/z^E_0$, with both
half marginals $\nu_E$ (\cref{lem:record:fold-identification}(d)).

\begin{lemma}[The exact-marginal matching law]\label{lem:record:e-law}
Let $p\ge19$ and $z_1=(p^2-3p+4)/(p(p-1))$.
\begin{enumerate}[label=(\alph*)]
\item The kernel
 $K_1=(1+\mathsf C-\delta_{\mathsf C})(1+\mathsf D-\delta_{\mathsf D})/z_1$ is the
 density of a $\mathcal G_p$-invariant probability measure on $H_E\times H_E$ whose
 two half marginals both equal $\nu_E$. It is $4/z_1$ times the indicator of the
 strict edges $x_1\to y_0$ and $x_0\to y_1$ and imposes no other relation between
 the halves; in particular other coincidences between the halves, such as
 $x_2=y_2$, have positive probability. Under it each of the pairs $(x_0,y_1)$ and
 $(x_1,y_0)$ is uniform on the strict-square pairs.
\item $z^E_0=(p-3)(p-4)(p-5)/(p(p-1)(p-2))\ge1-9q$, and on mean-zero functions
 $K^E_0$ has norm at most $k^E_{\rm old}=9q/(1-9q)$; its exact value there is
 $3/(p-3)$ (\cref{sec:frontier:empty-triple-negative}). Moreover $z_1\ge1-2q$ and
 $0<z_1-z^E_0=(7p^2-37p+52)/(p(p-1)(p-2))\le7q/Z$.
\item $K_1-K^E_0=\Pi_\perp(K_1-K^E_0)\Pi_\perp$, and
 \[
  \operatorname{Skew}(K_1-K^E_0)=\frac{\mathsf C+\mathsf D-\delta_{\mathsf C}\mathsf D
  -\delta_{\mathsf D}\mathsf C}{z_1},\qquad
  \|\mathsf C+\mathsf D\|=\frac r{1-q},\qquad
  \|\delta_{\mathsf C}\mathsf D\|,\|\delta_{\mathsf D}\mathsf C\|\le\frac{qr}{1-q},
 \]
 so that $\|\operatorname{Skew}(K_1-K^E_0)\|\le a^E_p:=r(1+2q)/((1-q)(1-2q))$.
\item Write $1-I_\times=\sum_{i,j}\delta_{ij}-\varepsilon_E$. Then $\varepsilon_E\ge0$,
 $\|\varepsilon_E\|\le18q^2/(1-q)$, and
 \[
  \Pi_\perp\operatorname{Sym}(K_1-K^E_0)\Pi_\perp
  =\mathcal S+\frac{\Pi_\perp\mathsf C\mathsf D\Pi_\perp
  +\Pi_\perp\delta_{\mathsf C}\delta_{\mathsf D}\Pi_\perp}{z_1}
  -\frac{\Pi_\perp\varepsilon_E\Pi_\perp}{z^E_0},
 \]
 where $\mathcal S=qV[J_3/z^E_0-(E_{10}+E_{01})/z_1]V^*$ on mean-zero functions,
 $E_{ij}$ being matrix units. Moreover $\|\mathcal S\|\le\mathfrak s_p$,
 $\|\mathsf C\mathsf D\|/z_1\le q/((1-q)(1-2q))$ and
 $\|\delta_{\mathsf C}\delta_{\mathsf D}\|/z_1\le q^2/((1-q)(1-2q))$, where
 \[
  \mathfrak s_p=\frac q{1-q}\Bigl[\frac{1+\sqrt2}{1-9q}
  +\frac{7q}{Z(1-9q)(1-2q)}\Bigr].
 \]
\end{enumerate}
\end{lemma}

\begin{proof}
(a) For $a\ne b$, $1+\chi(b-a)$ is $2$ if $a\to b$ and $0$ otherwise, and for
$a=b$ the factor $1+0-1$ vanishes; this describes $K_1$. Fix a half $x$. Under
$\nu_E$ the pair $(y_0,y_1)$ is uniform on ordered distinct pairs, and the pairs
with $x_1\to y_0$ and $x_0\to y_1$ number $(\frac{p-1}2)^2-k_p=(p^2-3p+4)/4$, the
coincident ones being the $k_p$ common out-neighbours of $x_0$ and $x_1$
(\cref{lem:paley-triple}(a)), whatever the orientation between $x_0$ and $x_1$.
Hence $\E_yK_1(x,y)=4((\frac{p-1}2)^2-k_p)/(p(p-1)z_1)=1$; the column sums are
one by the same count with common in-neighbours. Invariance holds because $\chi$
is invariant under square dilations; every relation other than the two strict
edges is free; and $\mathcal G_p$ acts transitively on the ordered pairs $(a,a')$
with $a\to a'$, so the invariant law of each queried pair is uniform on them.
(b) Given $x$, the half $y$ avoids three labels with probability $z^E_0$; the rest
is the symmetry of $I_\times$, the Schur test with row sums
$\E_y(1-I_\times)\le9q$, and the identities
$(1-3q)(1-4q)(1-5q)-(1-9q)Z=18q^2-42q^3$, $p^2-3p+4-(p-1)(p-2)=2$ and
$(p^2-3p+4)(p-2)-(p-3)(p-4)(p-5)=7p^2-37p+52\le7p^2$.
(c) Both laws have the half marginals $\nu_E$, which gives the first identity.
Expanding the product and using that transposition maps $\mathsf C,\mathsf D$
to $-\mathsf D,-\mathsf C$ and exchanges $\delta_{\mathsf C},\delta_{\mathsf D}$,
\[
 z_1\operatorname{Sym}K_1=1+\mathsf C\mathsf D-\delta_{\mathsf C}-\delta_{\mathsf D}
 +\delta_{\mathsf C}\delta_{\mathsf D},\qquad
 z_1\operatorname{Skew}K_1=\mathsf C+\mathsf D-\delta_{\mathsf C}\mathsf D
 -\delta_{\mathsf D}\mathsf C,
\]
and $K^E_0$ is symmetric. As operators $\mathsf C=V_1PV_0^*$ and
$\mathsf D=V_0PV_1^*$, so $\mathsf C+\mathsf D$ is
$\bigl(\begin{smallmatrix}0&1\\1&0\end{smallmatrix}\bigr)\otimes P$ conjugated by the
row $(V_0,V_1)$; the Gram matrix
$\bigl(\begin{smallmatrix}1&-u\\-u&1\end{smallmatrix}\bigr)$ of that row commutes with
$\bigl(\begin{smallmatrix}0&1\\1&0\end{smallmatrix}\bigr)$, with eigenvalues $1-u$ and
$1+u$ on $(1,1)$ and $(1,-1)$, so $\|\mathsf C+\mathsf D\|=(1+u)\|P\|=r/(1-q)$. The
kernel $\delta_{\mathsf C}\mathsf D$ depends on $(x_0,x_1)$ and $(y_0,y_1)$; for
independent labels it is $qI\otimes P$, of norm $qr$, and the two pair densities
$1/(1-q)$ give the bound; likewise for $\delta_{\mathsf D}\mathsf C$. Then use
$z_1\ge1-2q$.
(d) With $\kappa=\sum\delta_{ij}$, $\varepsilon_E=\kappa-\one_{\kappa\ge1}$ is bounded
by the sum of $\delta_{ij}\delta_{i'j'}$ over the $18$ compatible pairs, and given
$x$ two prescribed labels of $y$ at two prescribed positions have probability
$1/(p(p-1))$; the Schur test gives the bound. Since
$\Pi_\perp\delta_{ij}\Pi_\perp=qV_iV_j^*$ on mean-zero functions, the compressions of
$\sum\delta_{ij}/z^E_0$ and of $(\delta_{\mathsf C}+\delta_{\mathsf D})/z_1$ give
$\mathcal S$, and the constants $1/z_1$, $1/z^E_0$ are removed by the
compression. Write $J_3/z^E_0-(E_{10}+E_{01})/z_1=(J_3-E_{10}-E_{01})/z^E_0
+(E_{10}+E_{01})(1/z^E_0-1/z_1)$. The matrix $J_3-E_{10}-E_{01}$ has the
characteristic polynomial $\lambda^3-3\lambda^2+\lambda+1$, with roots $1$ and
$1\pm\sqrt2$, and $\|E_{10}+E_{01}\|=1$; with $\|V\|^2=1/(1-q)$ and (b) this gives
$\|\mathcal S\|\le\mathfrak s_p$. The kernel $\mathsf C\mathsf D$ is a two-edge
matching on the pairs $(x_0,x_1)$, $(y_0,y_1)$, of norm $q$ for independent
labels, and $\delta_{\mathsf C}\delta_{\mathsf D}$ is $qI\otimes qI$ there; the pair
densities give the last two bounds.
\end{proof}

\begin{lemma}[Orthogonal channels]\label{lem:record:e-orthogonal}
Let $p\ge19$.
\begin{enumerate}[label=(\alph*)]
\item The pair-star map $\mathcal T_E\colon L^2(\Fp)\to L^2(\nu_E)$,
 $\mathcal T_Eg(x)=\E_t\chi(t-x_1)\chi(t-x_0)g(t)$, has norm at most
 $(q(1-2q)/(1-q))^{1/2}\le r$ on mean-zero functions: for independent labels its
 norm there is exactly $(q(1-2q))^{1/2}$, at most $(q(1-q))^{1/2}$.
\item On mean-zero functions $\mathsf C\mathsf DV_0=\mathsf C\mathsf DV_1=-u\mathcal T_E$
 and $\mathsf C\mathsf DV_2=2\mathcal T_E/((p-1)(p-2))$. Hence, if $H$ is the
 orthogonal projection onto $V(H_0^3)$,
 \[
  \|\mathsf C\mathsf DH\|\le e_E:=\frac{qr}{((1-q)(1-3q))^{1/2}}
  \Bigl(2+\frac{4q^2}{(1-2q)^2}\Bigr)^{1/2}.
 \]
\item $\|\operatorname{Sym}(K_1-K^E_0)\|\le s^E_p$, where
 \[
  s^E_p=\mathfrak s_p+\frac{2e_E}{1-2q}+\frac{q^2}{(1-q)(1-2q)}
  +\frac{18q^2}{(1-q)(1-9q)} .
 \]
\end{enumerate}
\end{lemma}

\begin{proof}
(a) For independent labels $\mathcal T_E^*\mathcal T_E$ has the kernel
$(\E_x\chi(t-x)\chi(t'-x))^2=(\one_{t=t'}-q)^2$ by \cref{lem:paley-triple}(b), so
$\mathcal T_E^*\mathcal T_E=q(1-2q)I+q^2\Jproj$; restricting the pair $(x_0,x_1)$ to
distinct labels costs the density factor $1/(1-q)$.
(b) Given $y_0$, the label $y_1$ is uniform off $y_0$, so
$\E[\chi(y_1-x_0)\mid y_0]=-\chi(y_0-x_0)/(p-1)$; this gives
$\mathsf C\mathsf DV_0=-u\mathcal T_E$, and $\mathsf C\mathsf DV_1=-u\mathcal T_E$ in
the same way. Given $y_2=t$, the pair $(y_0,y_1)$ is uniform on the ordered
distinct pairs avoiding $t$, and by \cref{lem:paley-triple}(b)
$\sum_{a\ne a';\,a,a'\ne t}\chi(a-x_1)\chi(a'-x_0)=1+2\chi(t-x_1)\chi(t-x_0)$,
whose constant term is annihilated on mean-zero $g$; this gives
$\mathsf C\mathsf DV_2$. If $h=Vc$ with $c$ a triple of mean-zero functions, then
$\|h\|^2\ge(1-2u)\|c\|^2$ and
$\mathsf C\mathsf Dh=\mathcal T_E(-uc_0-uc_1+2c_2/((p-1)(p-2)))$, so
$\|\mathsf C\mathsf Dh\|\le ru(2+4/(p-2)^2)^{1/2}(1-2u)^{-1/2}\|h\|=e_E\|h\|$.
(c) By \cref{lem:record:e-law}(c),(d) it suffices to bound
$\|\mathcal S+X/z_1\|$ with $X=\Pi_\perp\mathsf C\mathsf D\Pi_\perp$. Here
$\mathcal S=H\mathcal SH$, and $X$ is self-adjoint with
$XH=\Pi_\perp\mathsf C\mathsf DH$, because $\Pi_\perp H=H$; so $\|XH\|=\|HX\|\le e_E$
and $\|X-(I-H)X(I-H)\|=\|HX+(I-H)XH\|\le2e_E$. The operators $\mathcal S$ and
$(I-H)X(I-H)/z_1$ act on orthogonal subspaces, so their sum has norm at most
$\max(\|\mathcal S\|,\|X\|/z_1)\le\mathfrak s_p$, because
$(1+\sqrt2)/(1-9q)\ge1/(1-2q)$. Hence $\|\mathcal S+X/z_1\|\le\mathfrak s_p+2e_E/z_1$,
and the other two terms of \cref{lem:record:e-law}(d) contribute
$q^2/((1-q)(1-2q))$ and $18q^2/((1-q)(1-9q))$.
\end{proof}

\begin{lemma}[Prices of the matching records]\label{lem:record:exact-marginal-e}
Let $p\ge2^{20}$. In the comparison law each two-layer matching record $E_0,E_1$
of a middle block is replaced, at every prime, by the law $K_1$ of
\cref{lem:record:e-law}, and the block record is the tensor product over the
primes.
\begin{enumerate}[label=(\roman*)]
\item On $L^2(\nu_E)$, $\|\operatorname{Sym}(K_1-K^E_0)\|\le s^E_p$ and
 $\|\operatorname{Skew}(K_1-K^E_0)\|\le a^E_p$. The centred operator of $K^E_0$ is
 self-adjoint, of norm $k^E_p=3/(p-3)=3q/(1-3q)$
 (\cref{sec:frontier:empty-triple-negative}), so the centred symmetric part of
 $K_1$ has norm at most $c^{\rm s}_Eq:=k^E_p+s^E_p$ and the centred operator of
 $K_1$ norm at most $\kappa^E_p:=k^E_p+s^E_p+a^E_p$. Each of $s^E_p/r$, $a^E_p/r$,
 $c^{\rm s}_E$ and $\kappa^E_p/r$ is a nondecreasing function of $(q,r)$.
\item Let $\Blk$ be a block of a piece starting at $X\ge2^{20}$, with least prime
 $b$, let $q_X=1/X$ and $r_X\ge X^{-1/2}$, and put, with every quantity of (i) at
 $(q_X,r_X)$,
 \[
 \begin{gathered}
  s_E(X)=\frac{s^E}r,\qquad a_E(X)=\frac{a^E}r=\frac{1+2q_X}{(1-q_X)(1-2q_X)},\\
  c^{\rm s}_E(X)=c^{\rm s}_E,\qquad\kappa_E(X)=\frac{\kappa^E}r,\qquad
  r_E(X)=2c^{\rm s}_E(X)a_E(X).
 \end{gathered}
 \]
 Suppose that, at $(q_X,r_X)$, $3q/(1-3q)<1$, $\kappa_Er<1$, $r_Er^2<1$,
 $4/((1-q)(1-2q))<r_E$, $\kappa_E^3<r_E$,
 \[
  s^E_{(2)}:=\frac{6r}{1-3q}s_E+s_E^2+a_E^2<2\qquad\text{and}\qquad
  s^E_{(3)}:=\Bigl[\Bigl(\frac{3r}{1-3q}\Bigr)^3+\kappa_E^3\Bigr]r<2 .
 \]
 Then the block differences
 $\bigotimes_{p\in\Blk}K_{1,p}-\bigotimes_{p\in\Blk}K^E_{0,p}$ of $E_0$ and $E_1$ have,
 on $L^2(\bigotimes_p\nu_{E,p})$, symmetric part of norm at most $s_E(X)b^{-1/2}$ and
 skew part of norm at most $a_E(X)b^{-1/2}$. The block law is invariant under
 independent maps of $\mathcal G_p$ at the primes of $\Blk$, accepts the six queries,
 and its queried pair is uniform on the pairs that are strict-square at every
 prime of $\Blk$.
\end{enumerate}
\end{lemma}

\begin{proof}
(i) The bounds are \cref{lem:record:e-law}(c) and \cref{lem:record:e-orthogonal}(c);
the centred operator of $K_1$ is that of $K^E_0$ plus the difference, and its
symmetric part is that of $K^E_0$ plus the symmetric part of the difference. The
monotonicity holds because every term is a positive expression in $(q,r)$ with
nonnegative coefficients and positive decreasing denominators, after the powers
of $r$ are cancelled ($\mathfrak s_p/q$, $e_E/q$ and the last two terms of $s^E_p$
divided by $q$ are such expressions).
(ii) For $p\in\Blk$, $q_p\le q_X$ and $r_p\le b^{-1/2}\le X^{-1/2}\le r_X$, so by (i)
$s^E_p\le s_E(X)r_p$, $a^E_p\le a_E(X)r_p$, the centred symmetric part of $K_{1,p}$
has norm at most $c^{\rm s}_E(X)q_p$, the centred operator of $K_{1,p}$ norm at most
$\kappa_E(X)r_p<1$, and the centred reference $T_{0,p}$, self-adjoint, norm
$k^E_p\le(3r_X/(1-3q_X))r_p<1$. The prime operators are Markov operators of
couplings with exact marginals $\nu_{E,p}$, so by \cref{lem:sector}(a),(b) the
block difference vanishes on the empty sector, is the prime difference on the
singleton sectors, and the norms of its symmetric and skew parts are the
maxima over the sectors. On a singleton $\{p\}$ they are at most
$s_E(X)b^{-1/2}$ and $a_E(X)b^{-1/2}$. On a sector $\{p,p'\}$,
\cref{lem:sector}(h) gives the symmetric part
$T^\circ_{0,p}\otimes S^\circ_{p'}+S^\circ_p\otimes T^\circ_{0,p'}+S^\circ_p\otimes S^\circ_{p'}
+A^\circ_p\otimes A^\circ_{p'}$, of norm at most $s^E_{(2)}r_pr_{p'}\le s^E_{(2)}/b$, and the
skew part $(T_{0,p}+S_p)^\circ\otimes A^\circ_{p'}+A^\circ_p\otimes(T_{0,p'}+S_{p'})^\circ$,
of norm at most $2c^{\rm s}_E(X)a_E(X)b^{-3/2}=r_E(X)b^{-3/2}$; here $S_p$, $A_p$ are
the symmetric and skew parts of the prime difference. On a sector $S$ with
$|S|\ge3$ the reference tensor is self-adjoint, so the skew part is that of the
new tensor (\cref{lem:sector}(b)), of norm at most $(\kappa_E(X)b^{-1/2})^3<
r_E(X)b^{-3/2}$; and the symmetric part has norm at most the sum of the norms of
the two tensors, $[(3r_X/(1-3q_X))^3+\kappa_E(X)^3]b^{-3/2}\le s^E_{(3)}/b$; the further
centred factors are contractions. Since $s^E_b\ge\mathfrak s_b\ge(1+\sqrt2)q_b$ and
$s^E_b\le s_E(X)r_b$, the bounds $s^E_{(2)}/b,s^E_{(3)}/b<2/b$ are below
$s_E(X)b^{-1/2}$; and $r_E(X)b^{-3/2}\le r_E(X)r_X^2b^{-1/2}<b^{-1/2}\le
a_E(X)b^{-1/2}$. Invariance, acceptance and the queried pairs hold primewise by
\cref{lem:record:e-law}(a).
\end{proof}

The values of $s_E(X)$, $a_E(X)$, $r_E(X)$ at the $33$ piece starts, with
$r_X=1/\lfloor\sqrt X\rfloor$, and the gates of (ii) are those of
\cref{cert:record:e-coefficients}; at $Q_3$ they are those of
\cref{cert:signed-local}.

\begin{lemma}[Matching one-arm stars]\label{lem:oct24:matching-star}
Let $\Blk$ be a block of a piece starting at $X\ge2^{20}$, with least prime $b$,
under the hypotheses of \cref{lem:record:exact-marginal-e}(ii), and let
$D^{\rm a}_{E,\Blk}$ be the skew part of the block difference of a matching record.
The matching coefficients of the piece are
\[
 s_E(X)=\frac{s^E}r,\qquad a_E(X)=\frac{1+2q}{(1-q)(1-2q)},\qquad
 r_E(X)=2\Bigl(\frac3{1-3q}+\frac{s^E}q\Bigr)a_E(X),
\]
at $(q,r)=(q_X,r_X)$, with $s^E=s^E_p$ of \cref{lem:record:e-orthogonal}(c) there.
Then $D^{\rm a}_{E,\Blk}=L_{E,\Blk}+R_{E,\Blk}$, where $L_{E,\Blk}$ equals
$L_{E,p}=\mathsf C+\mathsf D=V_1PV_0^*-V_0(V_1P)^*$ on the sector of exact support
$\{p\}$, for $p\in\Blk$, and zero on every other sector. The leading term is a
one-arm star with centre $\mathcal H_0$ and arm in $\mathcal H_1$, of arm norm at most
$b^{-1/2}$ and norm at most $b^{-1/2}/(1-b^{-1})$; and
$\|R_{E,\Blk}\|\le r_E(X)b^{-3/2}$, where $r_E(X)\ge2(4+\sqrt2)$, and $r_E(X)<11$ if
$r_X\le2^{-10}$, for instance for $r_X=1/\lfloor\sqrt X\rfloor$.
\end{lemma}

\begin{proof}
At a prime, \cref{lem:record:e-law}(c) gives
$D^{\rm a}_{E,p}-L_{E,p}=(z_1^{-1}-1)(\mathsf C+\mathsf D)-(\delta_{\mathsf C}\mathsf D+
\delta_{\mathsf D}\mathsf C)/z_1$, and $1-2q\le z_1\le1$, so
$\|D^{\rm a}_{E,p}-L_{E,p}\|\le\frac r{1-q}\cdot\frac{2q}{1-2q}+\frac{2qr}{(1-q)(1-2q)}
=\frac{4qr}{(1-q)(1-2q)}$, which is below $r_E(X)b^{-3/2}$ by the gate
$4/((1-q)(1-2q))<r_E$ and monotonicity. On the sectors of size at least two the
skew part of the block difference is at most $r_E(X)b^{-3/2}$ by the proof of
\cref{lem:record:exact-marginal-e}(ii), and $L_{E,\Blk}$ vanishes there. The
sectors are orthogonal, which gives the bound on $R_{E,\Blk}$. The arm $V_1P$ has
norm $r_p\le b^{-1/2}$, and $\|L_{E,p}\|=r_p/(1-q_p)\le b^{-1/2}/(1-b^{-1})$
(\cref{lem:record:e-law}(c)); $L_{E,\Blk}$ is the orthogonal sum of the $L_{E,p}$
over the singleton sectors. Finally $3/(1-3q)\ge3$, $s^E/q\ge\mathfrak s_p/q\ge
1+\sqrt2$ and $a_E\ge1$ give $r_E\ge2(4+\sqrt2)$; and $r_E$ is a nondecreasing
function of $(q,r)$ whose value at $(2^{-20},2^{-10})$ is
$10.834\ldots<11$ by outward evaluation of its formula, so
$r_E(X)<11$ when $r_X\le2^{-10}$.
\end{proof}

\subsection{Merges replaced at fixed marginals, and deleted records}\label{subsec:merge-repair}

The six two-layer merge records of a middle block have weight
$w_{\rm R}r^2$ each. We replace them by positive laws which accept the six
queries and preserve their actual half marginals. Four records use a pair
of mirror edges; the other two use one edge with the unused labels drawn
conditionally from their original half laws. The mirror construction cancels
the leading symmetric error. The two-layer constant records also admit an
accepting replacement, so only their folded counterparts are deleted.

\begin{definition}[Merge records and their designated edges]\label{def:merge-records}
Let $E_0=(0,1,1,2,2,0)$ and $E_1=(0,0,1,1,2,2)$ be the matching patterns and,
for $s\in\{0,1\}$ and $k\in\{0,1,2\}$, let
$G_{sk}(j)=\one[E_s(j)=k]$ be the merge patterns, with objects $0,1$
(\cref{def:patterns}). At a prime $p\equiv3\pmod4$, the two-layer merge
record $\pi_{sk}$ is uniform on injective four-label tuples
$(x_0,x_1,y_0,y_1)$, where $x_a$ and $y_a$ are the labels of object $a$
in layers zero and one. At the reserved cut both half marginals are the
uniform law $m_p$ on ordered distinct pairs. For this pattern
$G_{sk}$, every query of $\mathcal E_6$ reads the same designated
pair $x_{G_{sk}(0)}\to y_{G_{sk}(2)}$, determined by $G_{sk}$.

For $(s,k)\in\{(0,0),(0,1),(1,0),(1,1)\}$ replace $\pi_{sk}$ by the
uniform injective four-label law satisfying the two mirror edges
$x_1\to y_0$ and $x_0\to y_1$. For $(s,k)\in\{(0,2),(1,2)\}$ use instead
the law with density
\[
 \frac{2}{1-1/p}\one[x_0\to y_0]
 \quad\hbox{relative to }m_p\otimes m_p.
\]
The latter law allows coincidences between the two halves, but each half
remains an ordered distinct pair. Block replacements are products of
these laws over the primes of the block.
\end{definition}

\begin{table}[ht]
\centering\small
\begin{tabular}{@{}llll@{}}
\toprule
record & pattern & designated $\mathcal E_6$ edge & replacement\\
\midrule
$G_{00}$ & $(1,0,0,0,0,1)$ & $x_1\to y_0$ & both mirror edges\\
$G_{01}$ & $(0,1,1,0,0,0)$ & $x_0\to y_1$ & both mirror edges\\
$G_{02}$ & $(0,0,0,1,1,0)$ & $x_0\to y_0$ & first-coordinate edge\\
$G_{10}$ & $(1,1,0,0,0,0)$ & $x_1\to y_0$ & both mirror edges\\
$G_{11}$ & $(0,0,1,1,0,0)$ & $x_0\to y_1$ & both mirror edges\\
$G_{12}$ & $(0,0,0,0,1,1)$ & $x_0\to y_0$ & first-coordinate edge\\
\bottomrule
\end{tabular}
\caption{The six merge records and their accepting replacements. Each
replacement has exactly the original ordered-distinct-pair half laws.}
\label{tab:merge-records}
\end{table}

\begin{lemma}[Exact merge repairs]\label{lem:merge-repair}
The replacements of \cref{def:merge-records} are positive probability laws
with both actual half marginals equal to $m_p$. They are invariant under
$\mathcal G_p$, accept every query of $\mathcal E_6$, and have uniform
queried pairs on the ordered strict-square pairs.

Write $D_p$ for the old Markov operator on $L^2(m_p)$, and $P_p^{\rm c}$,
$P_p^{\rm s}$ for the mirror-edge and first-coordinate replacements.
For $p\ge2^{24}$ put
\begin{gather*}
 d_p=\frac2{p-2},\qquad \eta^{\rm G}_p=\frac1{p-1},\qquad
 \gamma^{\rm G}_p=\frac{\sqrt p}{p-1},\qquad
 \rho_p=\frac{\sqrt{p+1}}{p-1},\\
 s^{\rm G}_p=\frac{5p-6}{(p-2)(p-3)},\qquad
 a^{\rm G}_p=\frac{p^{-1/2}[1+2p^{-1/2}(4-6/p)]}{(1-2/p)(1-3/p)}.
\end{gather*}
For a block $\Blk$ with least prime $b$, the differences between the product
replacement and the product source have symmetric and skew norms at most
\[
 \begin{array}{c|cc}
 &\text{symmetric}&\text{skew}\\ \hline
 \text{mirror-edge replacement}&s_b&a_b\\
 \text{first-coordinate replacement}&d_b+\eta_b&\gamma_b.
 \end{array}
\]
All norms are relative to the product of the actual half laws $m_p$.
\end{lemma}

\begin{proof}
\emph{Half marginals and acceptance.}
Fix $x_0\ne x_1$ and put $k=(p-1)/2$. Among labels different from both
$x_0,x_1$, the two allowed out-neighbour sets for $y_0,y_1$ have sizes
$k$ and $k-1$ in some order. Their intersection has size $(p-3)/4$.
Thus the number of admissible distinct pairs $(y_0,y_1)$ is
\[
 k(k-1)-(p-3)/4=(p-2)(p-3)/4,
\]
independent of the fixed half. The incoming-neighbour count is identical,
so the mirror law has both marginals $m_p$. The first-coordinate law has
the same marginals because each coordinate of $m_p$ is uniform and a
strict Paley edge has uniform endpoints. Each designated edge in
\cref{tab:merge-records} holds under its replacement. Square-affine
invariance and transitivity on ordered strict-square pairs give the
claimed queried-pair law. The construction uses only $\mathcal E_6$.

\emph{The source and the first-coordinate law.}
The source operator averages over ordered pairs disjoint from its input
pair. Its eigenvalues are $1$ on constants, zero on antisymmetric pair
functions, $-2/(p-2)$ on $u(x_0)+u(x_1)$ with $\E u=0$, and
$2/((p-2)(p-3))$ on symmetric functions with zero row sums. These mutually
orthogonal spaces exhaust $L^2(m_p)$: summing a function over all pairs
and subtracting the pairs meeting the input gives the displayed values.
In particular $\|D_p-\Jproj\|=d_p$.

Let $V_i$ embed a singleton function as a function of coordinate $i$ of
$m_p$. The normalized strict Paley operator $P$ has symmetric part
$\Jproj-\eta^{\rm G}_p(I-\Jproj)$, skew norm $\gamma^{\rm G}_p$, and centered norm
$\rho_p$. Since $P_p^{\rm s}=V_0PV_0^*$, its difference from $D_p$ has
symmetric norm at most $d_p+\eta^{\rm G}_p$, skew norm $\gamma^{\rm G}_p$, and centered
replacement norm $\rho_p$.

\emph{The mirror law.}
Write $C=\chi(y_0-x_1)$ and $H=\chi(y_1-x_0)$. On the injective source,
$P_p^{\rm c}$ has density $(1+C)(1+H)$ relative to $D_p$.
Transposition exchanges $C,H$ with $-H,-C$, respectively. Hence, as
identities of kernels relative to $m_p\otimes m_p$,
\begin{equation}\label{eq:merge:mirror-identities}
 \operatorname{Sym}(P_p^{\rm c}-D_p)=D_pCH,\qquad
 \operatorname{Skew}(P_p^{\rm c}-D_p)=D_p(C+H),
\end{equation}
where the products on the right are pointwise products of densities.
Put $z=(p-2)(p-3)/(p(p-1))$ and $\tau=1-z$. Then $D_p$ is the
indicator of disjoint halves divided by $z$, and the deleted collision
kernel has constant row and column mass $\tau$. On independent field
pairs the two-edge character kernel $CH$ has norm $1/p$; restricting
both halves to distinct pairs costs at most $(1-1/p)^{-1}$. Schur's
bound on the collision kernel gives
\[
 \|D_pCH\|\le\frac{1/(p-1)+\tau}{z}
 =\frac{5p-6}{(p-2)(p-3)}=s^{\rm G}_p.
\]
For the linear term, the singleton embeddings on mean-zero functions
satisfy $V_i^*V_i=I$ and $V_0^*V_1=V_1^*V_0=-\eta^{\rm G}_p I$.
If $Kf(x)=\E_y\chi(y-x)f(y)$, then
$C+H=V_1KV_0^*+V_0KV_1^*$. Orthogonalizing $V_0+V_1$ and
$V_0-V_1$ gives the two blocks $(1-\eta^{\rm G}_p)K$ and
$-(1+\eta^{\rm G}_p)K$, so the norm is $(1+\eta^{\rm G}_p)/\sqrt p=\gamma^{\rm G}_p$.
The collision part has norm at most $2\tau$. Consequently
\[
 \|D_p(C+H)\|\le\frac{\gamma^{\rm G}_p+2\tau}{z}=a^{\rm G}_p,
 \qquad \|P_p^{\rm c}-\Jproj\|\le\alpha^{\rm G}_p:=d_p+s^{\rm G}_p+a^{\rm G}_p.
\]

\emph{Products over primes.}
Every operator is Markov in the same actual half law, so it preserves
the orthogonal decomposition by exact nonconstant prime support.
For $p\ge2^{24}$ the displayed formulas imply
\[
 \alpha^{\rm G}_p\le2p^{-1/2},\quad d_p\le p^{-1/2},\quad
 s^{\rm G}_p\ge5/p,\quad a^{\rm G}_p\ge4/p,
 \qquad \rho_p^2+d_p^2\le d_p+\eta^{\rm G}_p,\quad
 \rho_p^2\le\gamma^{\rm G}_p.
\]
For example, the first two follow by substituting
$r\le2^{-12}$ in their positive rational majorants; the other four
follow directly after clearing denominators. All the bounds decrease
with $p$. On singleton supports use the prime bounds. On a support of
size at least two, the mirror difference has symmetric norm at most
$\alpha_b^2+d_b^2\le5/b\le s_b$, and skew norm at most
$\alpha_b^2\le4/b\le a_b$, since the source is self-adjoint. Further
support factors are contractions. For the first-coordinate difference
the corresponding bounds are $\rho_b^2+d_b^2$ and $\rho_b^2$.
Taking the maximum over the orthogonal supports proves the block claims.
\end{proof}

\begin{corollary}[The six merge prices]\label{cor:merge-repair}
For a middle piece starting at $X$, set $q_X=1/X$,
$r_X=1/\lfloor\sqrt X\rfloor$, and
\begin{align*}
 g_{\rm S}(X)&=\frac{20}{(1-2q_X)(1-3q_X)}
                  +\frac4{1-2q_X}+\frac2{1-q_X},\\
 g_{\rm A}(X)&=\frac{4(1+8r_X)}{(1-2q_X)(1-3q_X)}
                  +\frac2{1-q_X}.
\end{align*}
On a block of that piece, with least prime $b$ and $r=b^{-1/2}$, the six
merge changes have total symmetric and skew prices at most
\[
 w_{\rm R}g_{\rm S}(X)r^4q_\Blk^{2a},\qquad
 w_{\rm R}g_{\rm A}(X)r^3q_\Blk^{2a},
\]
respectively. Their norm envelopes are positive reserved diagonals of
their actual ordered-distinct-pair half laws.
\end{corollary}

\begin{proof}
There are four mirror and two first-coordinate records, each of mass
$w_{\rm R}r^2$ before the common normalization. In the formulas of
\cref{lem:merge-repair}, drop the negative numerator corrections and
use $1/b\le q_X$, $b^{-1/2}\le r_X$. Summing gives the displayed
coefficients. The charging rule of \cref{rem:cp:charging}(ii) applies
relative to the unchanged actual half marginals. Omitting the
division by the reference normalizer (a number $\ge1$, so its
reciprocal is $\le1$) only enlarges the prices.
\end{proof}

\begin{certificate}[The merge coefficients]\label{cert:merge-coefficients}
The coefficients $g_{\rm S}(X)$ and $g_{\rm A}(X)$ are the explicit positive
rational functions of \cref{cor:merge-repair}; no graph-count approximation
is needed to define them. At $X=Q_3=19951586$ they are, respectively,
$26.0000055133\ldots$ and $6.0071663530\ldots$, and they decrease to $26$
and $6$. The exact checker
\path{research/four_routes_2026_10_26/pair_transfer/actual_g_mirror_check.py}
verifies the rational endpoint gates at $p=2^{24}$ and the physical
marginals, designated edges and mirror identities by exact enumeration
at $p=7,11,19$. It also checks the source spectral polynomial at $p=7,11$.
The uniform proof and the operator identities were independently reviewed
in \path{research/four_routes_2026_10_26/folded_states/INDEPENDENT_ACTUAL_G_MIRROR_REVIEW.md}.
The original source weights and reference laws are unchanged; these
coefficients replace only the merge comparison prices.
\end{certificate}

\begin{lemma}[Accepting two-layer constants]\label{lem:record:accepting-c2}
Each of the two two-layer constant records of a block can be replaced
by a positive accepting law with the same actual half marginals. If
$b\ge7$ is the least prime (as for every middle block) and $r=b^{-1/2}$,
the symmetric and skew norms of
the normalized replacement minus the source are at most
\[
 \frac{r^2}{(1-1/b)^2},\qquad \frac r{1-1/b},
\]
respectively. For a piece starting at $X$, the two records together have
symmetric and skew prices at most
\[
 \frac{2w_{\rm C}}{(1-1/X)^2}r^5q_\Blk^{2a},\qquad
 \frac{2w_{\rm C}}{1-1/X}r^4q_\Blk^{2a}.
\]
Both norm envelopes are uniform one-label reserved diagonals.
\end{lemma}

\begin{proof}
Write $(a_p,b_p)$ for the labels of the single object of either
constant record in layers $0$ and $1$ at the prime $p$ (the labels
seen at the reserved cut). They are uniform
ordered distinct labels: this is the one-object marginal of the original
injective matching law. All six queries read $a_p\to b_p$. Replace this
law by the uniform strict Paley edge, independently at every prime.
It is square-affine invariant, accepts all six queries and has the same
uniform endpoint marginals. The two records give the same physical law,
but both of their weights are retained.

On uniform field $L^2$, the old and new operators are
$D_p=\Jproj-\eta_p(I-\Jproj)$ and
$P_p=D_p+\frac p{p-1}K_p$, where
$K_pf(a)=\E_b\chi(b-a)f(b)$, $K_p^*=-K_p$ and
$K_p^*K_p=p^{-1}(I-\Jproj)$. They commute and are normal. On a
nonconstant complex Fourier mode their eigenvalues are $-\eta_p$ and
$-\eta_p\pm i\gamma_p$. The singleton-support difference is therefore
purely skew, of norm $\gamma_p$. On a support of exactly two primes the
symmetric difference is $\pm\gamma_p\gamma_{p'}$. On a support of size
at least three its symmetric norm is at most
$\rho_b^3+\eta_b^3\le\gamma_b^2$; the last inequality follows from
$(b+1)^{3/2}+1\le b(b-1)$ for $b\ge7$.
On every support of size at least two the skew norm is at most
$\rho_b^2\le\gamma_b$. Constants cancel. These exact-support bounds
prove the asserted whole-block norms without a factor for the number
of primes. Each original record has mass $w_{\rm C}r^3$ before normalization,
so the two norm prices follow from \cref{rem:cp:charging}(ii).
The comparison replacement changes neither the reference nor its first
stage or Gram bracket.
\end{proof}

\paragraph{Deletions.}
The comparison law deletes the six folded merge records
($6v_{\rm R}r^3$), the polarization correction $G$ of the first stage
($2v_{\rm D}(f_1+h_1)r^3$), and the two folded constant records
($2v_{\rm C}r^4$). These are positive records, symmetric at $\iota$
(the folded ones reserved diagonal), and a deleted positive record is
charged its mass to the symmetric part only
(\cref{rem:cp:charging}(i)). The two-layer constant records are replaced
as in \cref{lem:record:accepting-c2} and are not included in these deletions.

\subsection{The comparison law of a middle block and its prices}
\label{subsec:middle-comparison}

On a middle block $\Blk$ the comparison law is
\[
 \nu_\Blk=\lambda_\Blk+\sum_{t\in\mathcal T_\Blk}\zeta_t,
\]
where $\zeta_t$ is the replacement of the record $t$ minus $t$, or minus $t$ if
$t$ is deleted, each with its weight in $\lambda_\Blk$, divided
by the normalizing mass $\widetilde\mu_\Blk(1)$ of the reference
(\cref{lem:front12:reference}). The records $t\in\mathcal T_\Blk$ and their
replacements are: the two two-layer fold records $F_1,F_2$ (weight
$w_{\rm D}r$), replaced by the signed records of
\cref{def:record:mixture}(c),(d); the two two-layer matching records (weight
$w_{\rm S}r$), replaced by the block laws of \cref{lem:record:exact-marginal-e};
the six two-layer merge records (weight $w_{\rm R}r^2$), replaced by the repairs
of \cref{cor:merge-repair}; the two two-layer constant records (weight
$w_{\rm C}r^3$), replaced by \cref{lem:record:accepting-c2}; and the deleted records listed above. The first
stage changes only one-layer folded records (\cref{def:record:frame}), so the
two-layer records enter $\lambda_\Blk$ with the weights of $\mu_\Blk$. The
retained records of $\nu_\Blk$ are the two-layer root and $F_0$, the folded root
and the folded $F_0$, the record $E_{\rm new}$ of the first stage, and the
replacements; each accepts the six queries: the root and $F_0$ contain the edge
$0\to2$ of the first coordinate in both layers, and so do their diagonal copies
(on the pieces from $2^{30}$ the root is the positive root of
\cref{sec:positive-roots}, which has the support of the uniform root and
accepts the six queries, \cref{lem:positive-root-properties,imp:local:filtered-coherent,cert:positive-root-caps}); $E_{\rm new}$ by
\cref{lem:record:one-stage-folded-e}; the replacements by
\cref{lem:signed-components,lem:record:exact-marginal-e,cor:merge-repair,lem:record:accepting-c2}. The
comparison law of the block is a signed measure,
\[
 \nu_\Blk=\nu^+_\Blk-\frac\tlaw2\,w\,\Pi^{--}_{F_1,\Blk}-\frac\tlaw2\,w\,\Pi^{--}_{F_2,\Blk},
 \qquad w=\frac{w_{\rm D}r}{\widetilde\mu_\Blk(1)},
\]
where $\Pi^{--}_{F_1,\Blk}$ and $\Pi^{--}_{F_2,\Blk}$ are the probability records
placed by $F_1$ and $F_2$ with the laws $\bigotimes_pP_{--,p}$ and
$\bigotimes_pP^\varsigma_{--,p}$, and $\nu^+_\Blk\ge0$ consists of all other records
of $\nu_\Blk$, among them the components $\bigotimes_pP_{++,p}$ and
$\frac\tlaw2\bigotimes_pP_{-+,p}$ of the two signed records.

\begin{lemma}[Comparison prices of a middle block]\label{lem:middle-comparison-prices}
Let $\Blk$ be a block of the middle piece $i$ starting at $X$, with least prime
$b\ge X$ and $r=b^{-1/2}$, and evaluate the constants of
\eqref{eq:signed-constants} and the matching coefficients of
\cref{lem:record:exact-marginal-e}(ii) at the piece start, at $(q_X,r_X)$ with
$q_X=1/X$ and $r_X=1/\lfloor\sqrt X\rfloor$. Then
$\nu_\Blk$ is a finite real combination of the three nonnegative measures
$\nu^+_\Blk$, $\Pi^{--}_{F_1,\Blk}$ and $\Pi^{--}_{F_2,\Blk}$, each invariant under
independent square-affine maps at the primes of $\Blk$ and accepting the six
queries, with $\nu_\Blk(1)>0$ and total variation
$|\nu|_\Blk\le\nu_\Blk(1)+2\tlaw w$; it satisfies (a)
of \cref{def:record:comparison-law}, and its changes satisfy (b) there with norm
envelopes $\operatorname{Diag}_\iota$ of their actual half marginals and prices
\[
 s_\Blk\le\bigl(S_2r^2+S_3r^3+S_4r^4+S_5r^5\bigr)q_\Blk^{2a},\qquad
 a_\Blk\le\bigl(A_2r^2+A_3r^3+A_4r^4\bigr)q_\Blk^{2a},
\]
where
\begin{align*}
 S_2&=2w_{\rm S}s_E,&
 A_2&=2w_{\rm D}\ell_\tlaw+2w_{\rm S}a_E,\\
 S_3&=2w_{\rm D}s_F+6v_{\rm R}+2v_{\rm D}(f_1+h_1),&
 A_3&=2w_{\rm D}a_F+w_{\rm R}g_{\rm A},\\
 S_4&=2v_{\rm C}+w_{\rm R}g_{\rm S},&
 A_4&=\frac{2w_{\rm C}}{1-q_X},\\
 S_5&=\frac{2w_{\rm C}}{(1-q_X)^2}.&&
\end{align*}
with $s_F,a_F,\ell_\tlaw$ of \eqref{eq:signed-constants} (\cref{prop:signed-bounds}),
$s_E=s_E(X)$, $a_E=a_E(X)$ of \cref{lem:record:exact-marginal-e}(ii), and
$g_{\rm S}=g_{\rm S}(X)$, $g_{\rm A}=g_{\rm A}(X)$ of \cref{cor:merge-repair}. The auxiliary prices
(\cref{def:cp5}) of the skew change of each of the two fold records are, times
$q_\Blk^{2a}$, the raw-arm price $w_{\rm D}r^2$, the leading price
$w_{\rm D}\ell_\tlaw r^2$ and the remainder price $w_{\rm D}a_Fr^3$; those of each
of the two matching records are $w_{\rm S}r^2$, $w_{\rm S}r^2/(1-q_X)$ and
$w_{\rm S}r_Er^4$, with $r_E=r_E(X)$ of \cref{lem:oct24:matching-star}. Summed over
the two records of each kind they are $2w_{\rm D}r^2$, $2w_{\rm D}\ell_\tlaw r^2$,
$2w_{\rm D}a_Fr^3$ and $2w_{\rm S}r^2$, $2w_{\rm S}r^2/(1-q_X)$, $2w_{\rm S}r_Er^4$.
Every selected law of these changes is admissible (\cref{def:admissible}).
\end{lemma}

\begin{proof}
Dividing by $\widetilde\mu_\Blk(1)\ge1$ only decreases the weights, so it suffices to
charge the unnormalized ones. \emph{Decomposition and (a).} The records of
$\lambda_\Blk$ that the comparison law keeps, the replacements of the matching
and merge records, the accepting two-layer constants, and the components $\bigotimes_pP_{++,p}$, $\frac\tlaw2\bigotimes_pP_{-+,p}$
of the two signed records are nonnegative measures; their sum is $\nu^+_\Blk$,
and $\nu_\Blk(1)=\nu^+_\Blk(1)-\tlaw w$, which gives the total variation. The
two-layer root has weight one in $\widetilde\mu_\Blk$ and is retained, so
$\nu^+_\Blk(1)\ge1/\widetilde\mu_\Blk(1)$ and
$\nu_\Blk(1)\ge(1-\tlaw w_{\rm D}r)/\widetilde\mu_\Blk(1)>0$, as $w_{\rm D}<6$
(\cref{rem:middle-weights}, \cref{tab:data:middle-weights}) and $r<2^{-12}$. Every
record and component is invariant under independent square-affine maps at the
primes of $\Blk$ (\cref{lem:signed-components}(a),(d), \cref{lem:record:e-law}(a),
\cref{cor:merge-repair,lem:record:accepting-c2}; for the unchanged records \cref{lem:front12:reference}
and \cref{lem:record:one-stage-folded-e}), so its one-slot marginals are uniform
and its queried pairs are its mass times the uniform law on strict-square pairs; and each accepts
the six queries, as listed above. This is (a) and the clause \ref{it:record:R5}
for each of the three components. \emph{Prices.} Each replacement has the same
actual half marginals on both halves of $\iota$ as the record it replaces, so
\cref{rem:cp:charging}(ii) charges it its weight times the norms of the symmetric
and skew parts of the difference; for the signed records the proof of (ii) uses
only that the new record has mass one and the half marginals of the old one,
not that it is nonnegative. The two fold records cost at most
$w_{\rm D}r\cdot s_Fr^2$ and $w_{\rm D}r\cdot(\ell_\tlaw r+a_Fr^2)$ each
(\cref{prop:signed-bounds}(a),(b),(e)), the two matching records
$w_{\rm S}r\cdot s_Er$ and $w_{\rm S}r\cdot a_Er$ each
(\cref{lem:record:exact-marginal-e}(ii)), and the six merge records together
$w_{\rm R}g_{\rm S}r^4$ to the symmetric part and $w_{\rm R}g_{\rm A}r^3$ to the skew
part (\cref{cor:merge-repair}). The two-layer constants add $S_5r^5$ and
$A_4r^4$, respectively (\cref{lem:record:accepting-c2}). The deleted records are
charged their masses to the symmetric part (\cref{rem:cp:charging}(i)). The skew
change of a fold record splits into its leading three-arm star, of norm at most
$\ell_\tlaw r$ with arms of norm at most $r$ (\cref{prop:signed-bounds}(b),
\cref{lem:actual-features}(e)), and a remainder of norm at most $a_Fr^2$; that of
a matching record into its one-arm star, of norm at most $r/(1-q_X)$ with arm
norm at most $r$, and a remainder of norm at most $r_Er^3$
(\cref{lem:oct24:matching-star}); times the weights these are the auxiliary
prices. The coefficients are evaluated at the piece start and bound the norms at
every block of the piece (\cref{lem:majorant}, the constants being
nondecreasing in $(q,r)$), and the weights and the row constants $f_1,h_1$ are
those of the piece (\cref{lem:record:local-domain}(iii)). \emph{Admissibility.}
The changes are invariant signed measures on the field labels only, independent
of the indices, so (O1) holds; and the only operator bounds applied to them are
their own norms at $\iota$, with no shared label, so (O2) holds.
\end{proof}

Summing over the blocks of each piece with the charges of
\cref{lem:record:block-charge} and over the pieces gives the middle parts of $S$
and $A$ and the totals of the auxiliary prices directly. In the piece $i$, with
start $X_i$, weights $w_{{\rm D},i}=w_{{\rm S},i}$, the coefficients of
\cref{lem:middle-comparison-prices} at $X_i$, and the certified degree charges
$\bar I_{d,i}\ge\sum_\Blk b_\Blk^{-d/2}q_\Blk^{2a_0}$ of \cref{tab:record:middle-pieces},
the prices at the anchor $a_0$ are
\begin{equation}\label{eq:middle-piece-prices}
\begin{aligned}
 &A^{\rm arm}_{F,i}=\operatorname{up}(2w_{{\rm D},i}\bar I_{2,i}),\quad
  A^{\rm lead}_{F,i}=\operatorname{up}(2w_{{\rm D},i}\ell_\tlaw(X_i)\bar I_{2,i}),\quad
  A^{\rm rem}_{F,i}=\operatorname{up}(2w_{{\rm D},i}a_F(X_i)\bar I_{3,i}),\\
 &S_{F,i}=\operatorname{up}(2w_{{\rm D},i}s_F(X_i)\bar I_{3,i}),\qquad
  A_{F,i}=A^{\rm lead}_{F,i}+A^{\rm rem}_{F,i},\\
 &A^{\rm arm}_{E,i}=\operatorname{up}(2w_{{\rm S},i}\bar I_{2,i}),\quad
  A^{\rm lead}_{E,i}=\operatorname{up}\Bigl(\frac{2w_{{\rm S},i}\bar I_{2,i}}{1-q_{X_i}}\Bigr),\quad
  A^{\rm rem}_{E,i}=\operatorname{up}(2w_{{\rm S},i}r_E(X_i)\bar I_{4,i}),\\
 &A_{E,i}=2\operatorname{up}(w_{{\rm S},i}a_E(X_i)\bar I_{2,i}),\qquad
  S_{E,i}=\operatorname{up}(2w_{{\rm S},i}s_E(X_i)\bar I_{2,i}),\\
 &S_{{\rm oth},i}=\sigma_{s,i}+S^{\rm oth}_{3,i}\bar I_{3,i}
       +S_{4,i}\bar I_{4,i}+S_{5,i}\bar I_{3,i}/X_i,\\
 &A_{{\rm oth},i}=\sigma_{a,i}+w_{{\rm R},i}g_{\rm A}(X_i)\bar I_{3,i}
       +A_{4,i}\bar I_{4,i},
\end{aligned}
\end{equation}
with $\operatorname{up}(x)=10^{-30}\lceil10^{30}x\rceil$ and
$S^{\rm oth}_3=S_3-2w_{\rm D}s_F=6v_{\rm R}+2v_{\rm D}(f_1+h_1)$.
Here $\sigma_{s,i},\sigma_{a,i}\in[0,10^{-30})$ are the rounding
slacks of the corresponding merge and two-layer-constant rows of
the authenticated middle-price source named in
\cref{def:closing-quantities}(3): each is its stored upward
rounding minus its exact unrounded expression. Keeping these
slacks makes the update an exact rational replacement of the
two-layer constant and merge contributions as stored there. The degree-five charge needs no new prime enumeration,
since $b^{-5/2}\le X_i^{-1}b^{-3/2}$.
The totals at the exponent $a$ are $t$ times the sums over the $33$ pieces,
$t=e^{2(a-a_0)L}$ with $L=570$ (\cref{lem:record:repricing},
\cref{def:closing-quantities}): $A^{\rm arm}_F=t\sum_iA^{\rm arm}_{F,i}$ and so on for
$A^{\rm lead}_F$, $A^{\rm rem}_F$, $A_F$, $S_F$, $A^{\rm arm}_E$, $A^{\rm lead}_E$,
$A^{\rm rem}_E$, $A_E$, $S_E$, $S_{\rm oth}$ and $A_{\rm oth}$, with the middle totals
$S_{\rm mid}=S_F+S_E+S_{\rm oth}$ and $A_{\rm mid}=A_F+A_E+A_{\rm oth}$. The merge and constant replacements have different symmetric and skew
prices, as displayed in \eqref{eq:middle-piece-prices}. The full skew
price $A_E$ of the matching records is bounded directly; $A_F$ is the sum of the
leading and remainder prices, and $A^{\rm lead}_F\ge A^{\rm arm}_F$,
$A_E\ge A^{\rm arm}_E$ (\cref{cert:signed-local}). At the density-linear
baseline exponent $a=\PushA$, the following conservative upper totals
remain valid: the generic symmetric and skew prices
$S_{{\rm oth},i}$ and $A_{{\rm oth},i}$ in \eqref{eq:middle-piece-prices} are each no larger than the corresponding generic rows of
\cref{def:closing-quantities}(3). The totals at $a$, which are the
baseline totals used in \cref{sec:closing}, are
$A^{\rm arm}_F=A^{\rm arm}_E<0.42332183839047$,
$A^{\rm lead}_F<0.53897233394172$,
$A^{\rm rem}_F<0.00005658950959$,
$A_F<0.53902892345131$,
$A^{\rm lead}_E<0.42332183982722$,
$A^{\rm rem}_E<1.436778E-8$,
$A_E<0.42332184557419$,
$S_F<0.00013191008864$,
$S_E<0.00003092662078$,
$S_{\rm oth}<0.00052581574396$,
$A_{\rm oth}<0.00023870593198$,
$S_{\rm mid}<0.00068865245337$ and
$A_{\rm mid}<0.96258947495747$
(\path{research/push_2026_10_25/artifacts/author.json}, key
\texttt{raw\_totals}; \cref{tab:data:prices}). The final construction
uses the smaller prices in \eqref{eq:middle-piece-prices} and also
retains the common reference normalization; its complete new totals
are given in \cref{sec:four-routes-closing}.

\paragraph{Consequences and what to read next.}
From Part~II: \cref{lem:majorant} (evaluation at the piece start),
\cref{lem:sector} (sectors of block differences, Markov tensor maximum, the
closing clause of (b) and the two-support step (h)) and
\cref{rem:signed-mixture-sectors} (signed sums of products),
\cref{def:admissible}, and the Weil bound \cref{thm:weil}. From Part~III: the
signed comparison laws of \cref{def:cp4}, the auxiliary prices of
\cref{def:cp5}, the charging rules \cref{rem:cp:charging}, the energy toolkit
\cref{lem:energy-toolkit} (for \cref{lem:record:old-f-marginal}(d)),
\cref{lem:record:repricing}. From Part~IV:
\cref{lem:record:fold-identification} (the records at $\iota$ and their half
marginals),
\cref{lem:paley-triple} (Paley counts, the Paley operator, directed triples and
their features), \cref{lem:spec:affpair} (invariant functions),
\cref{lem:extension,lem:graphs:sequential} (count errors),
\cref{sec:frontier:empty-triple-negative} (the disjoint-triple operator),
\cref{prop:middle-negative,cor:influence-schur} (correlations of the merge graph
laws), \cref{lem:record:actual-correlation,lem:record:projection} (actual marginals,
projected kernels), and the repair package
\cref{lem:tr:repair,cor:markov-completion}.
From \cref{sec:middle-reference,sec:middle-frame}: the reference
(\cref{lem:front12:reference}) and its weights (\cref{rem:middle-weights}), the
frame (\cref{def:record:frame}), the block charges, the local domain and the
recorded F21 caps (\cref{lem:layer12-caps}); from \cref{sec:positive-roots}, the support of the
positive root. Data and certificates:
\cref{cert:record:middle-geometry,cert:signed-local,cert:merge-coefficients,cert:record:old-f-cap,cert:record:e-coefficients};
the totals of \cref{tab:data:prices}.

\paragraph{Hypotheses verified in this section (\cref{tab:hypotheses}).} (CP4) for the middle atoms, that is,
the signed decomposition of $\nu_\Blk$ into nonnegative accepting components with
uniform one-slot marginals, its variation and \ref{it:record:R5}, and the
brackets of every change with the prices and auxiliary prices of
\cref{lem:middle-comparison-prices}; the negative cap $k_{\rm D}=k_{\rm D}(X,L)$ used
in (CP1) for the middle atoms (\cref{lem:old-f-law}(e)); the centred norm
$\kappa_F$ used by the variance credits (\cref{lem:record:old-f-marginal}); the
direct-Gram coefficient $d_\tlaw$ (\cref{prop:signed-bounds}(d)); and the
admissibility clause (CP8)(a) for the changes of the middle records.
